\section{The hierarchy \texorpdfstring{$\HBtp$}{HB*tau+}}\label{s2:sec}

This section defines the hierarchy $\HBtp$ on which the whole construction
runs, and proves the structural facts about it that later sections use. The
hierarchy follows the round structure of the proof of \cite[Theorem~2]{BM}:
in each round, long cycles are removed, the remaining graph is cut by the
recursion of \cite[Lemma~14]{BM} with $s=0$ into pieces, and each large piece
is cut again into robust expanders. There are two changes inside a round.
First, the robust recursion is replaced by Lemma~14$^\tau$ below: at every
split, vertices with many deleted edges across the cut are moved into the
separator (the \emph{$\tau$-rules}), which caps the number of deleted edges
from any vertex into the non-duplicated vertices of a leaf (the thin-cut
lemma). Second, a \emph{guest cap}
(rule~(GC)) is added to the definition of light pre-parts. The leaves of this
recursion are still certified expanders, and the overlap between leaves stays
small (Lemma OV).

All constants in this section are the constants of $\HBtp$ (\RTb{I8}). The
hierarchy without the $\tau$-rules and without (GC) (``plain $\HBs$''), and the
hierarchy with the $\tau$-rules but without (GC) (``$\HBt$''), appear only in
the final remark of this section, for comparison. Round parameters are written
$\tHB_l$, $\tau_r$ and $\thGC_r$ (\RTa{M10}).

\emph{Standing assumption.} Throughout this section, as everywhere in this
text, $\Dstar$ denotes the constant of Definition~\ref{s1:defConstants}(iii):
it satisfies \Gc{1}--\Gc{4} (Condition~\ref{s1:condGamma}). The proofs use
these conditions either directly or through the results they apply; for
instance the termination part of Proposition~\ref{s2:propExists} uses
Proposition~\ref{s2:propDegRec}, and hence \Gc{1}.

\emph{Local notation.} In the first two subsections of this section the
letter $H$ denotes a node graph with $m=|H|$ vertices, and $(U,F)$, $N$,
$F_0$, $F_1$, $F''$, $\Hout$, $\Hin$, $U'$, $N'$, $N''$, $G_1$, $G_2$, $n_1$,
$n_2$, $\nu^*$ are local to these subsections, as are $S$ (a total leaf
size), $M$ (a size threshold) and $c$ (a constant) in
Lemma OV. Throughout, $\eps=2^{-5}$, $\log=\log_2$, and
$|H|\defeq|V(H)|$ (Convention~\ref{s1:convGraphs}).

\subsection{Witnesses, the \texorpdfstring{$\tau$}{tau}-rules and separation}

\begin{definition}[witness and its minimal part]\label{s2:defWitness}
Let $s\ge0$ and let $H$ be a graph on $m=|H|$ vertices that is not an
$(\eps,s)$-expander (Cited result~\ref{s1:citDef11}). Then $m\ge2$, since a
graph on one vertex has no set $U$ with $1\le|U|\le2m/3$. A \emph{witness} at
$H$ is a pair $(U,F)$ with $U\subseteq V(H)$ and $F\subseteq E(H)$ such that
\[
1\le|U|\le\tfrac23m,\qquad |F|\le s|U|,\qquad
|\Nbr_{H-F}(U)|<\frac{\eps|U|}{\log^2m}.
\]
Since $H$ is not an $(\eps,s)$-expander, a witness exists. For a witness
$(U,F)$ put
\[
N\defeq\Nbr_{H-F}(U),\qquad F_0\defeq E_H\bigl(U,\,V(H)\setminus(U\cup
N)\bigr).
\]
\emph{Fact.} $F_0\subseteq F$ and $\Nbr_{H-F_0}(U)=N$. In particular
$(U,F_0)$ is again a witness at $H$, with the same set $N$.
\deps{\ref{s1:citDef11}, \ref{s1:defConstants}}
\end{definition}

\begin{proof}[Proof of the Fact]
Let $uv\in F_0$ with $u\in U$ and $v\notin U\cup N$. If $uv\notin F$, then
$uv$ is an edge of $H-F$, so $v\in\Nbr_{H-F}(U)=N$, a contradiction. Hence
$F_0\subseteq F$. Consequently $H-F\subseteq H-F_0$, so
$N\subseteq\Nbr_{H-F_0}(U)$. Conversely, let $v\notin U$ have a neighbour
$u\in U$ with $uv\notin F_0$. If $v\notin N$ then $uv\in
E_H(U,V(H)\setminus(U\cup N))=F_0$, a contradiction; so $v\in N$. Thus
$\Nbr_{H-F_0}(U)=N$. Finally $|F_0|\le|F|\le s|U|$, so $(U,F_0)$ is a witness.
\end{proof}

\begin{definition}[the $\tau$-rules and the split]\label{s2:defTauRules}
Let $H$, $m$ and $s$ be as in Definition~\ref{s2:defWitness}, let $(U,F)$ be a
witness at $H$ with $N$ and $F_0$ as there, write $V\defeq V(H)$, and let
$\tau>s$ be real. The \emph{$\tau$-rules} are:
\begin{enumerate}
\item[(1)] $\Hout\defeq\{v\in U:\deg_{F_0}(v)\ge\tau\}$,\quad
$U'\defeq U\setminus\Hout$,\quad $N'\defeq N\cup\Hout$;
\item[(2)] $F_1\defeq E_H\bigl(U',V\setminus(U'\cup N')\bigr)$,\quad
$\Hin\defeq\{x\in V\setminus(U'\cup N'):\deg_{F_1}(x)\ge\tau\}$,\quad
$N''\defeq N'\cup\Hin$,\quad $F''\defeq E_H\bigl(U',V\setminus(U'\cup
N'')\bigr)$.
\end{enumerate}
The \emph{split of $H$ by the $\tau$-rules} has the two children
$G_1\defeq H[U'\cup N'']-F''$ and $G_2\defeq H-U'-E(G_1)-F''$, and its
\emph{deleted set} is $F''$. Every edge of $F''$ has one end in $U'$ and the
other outside $U'\cup N''$, so no edge of $F''$ lies inside $U'\cup N''$, and
every edge of $F''$ meets $U'$. Hence
\begin{equation}\label{s2:eqSplit}
G_1=H[U'\cup N''],\qquad G_2=H[V\setminus U']-E\bigl(H[N'']\bigr).
\end{equation}
\emph{Facts.}
\begin{enumerate}
\item[(a)] $U'\cap N''=\emptyset$, $U'\cup N'=U\cup N$, and
$F''\subseteq F_1\subseteq F_0\subseteq F$.
\item[(b)] $E(H)=E(G_1)\sqcup E(G_2)\sqcup F''$ (disjoint union),
$V(G_1)=U'\cup N''$ and $V(G_2)=V\setminus U'$. A vertex of $U'$ lies only in
$G_1$, a vertex of $N''$ lies in both children, and a vertex of
$V\setminus(U'\cup N'')$ lies only in $G_2$.
\item[(c)] (\RTb{I7}) Given the pair $(U,N)$, the objects $\Hout$, $U'$, $N'$,
$F_1$, $\Hin$, $N''$, $F''$, $G_1$, $G_2$ depend only on $(U,N)$, because
they are defined from $F_0=E_H(U,V\setminus(U\cup N))$. The choice of $F$
matters only through $N=\Nbr_{H-F}(U)$: every witness $(U,F)$ yields the
same split as its minimal part $(U,F_0)$. (It is not true that the choice of
$F$ is irrelevant, since $F$ determines $N$.)
\end{enumerate}
\deps{\ref{s2:defWitness}}
\fixes{\RTb{I7}}
\end{definition}

\begin{proof}[Proof of the Facts]
(a) $U'\subseteq U$ and $N\cap U=\emptyset$, so $U'\cap N=\emptyset$;
$\Hout\cap U'=\emptyset$ by definition; and $\Hin\cap(U'\cup N')=\emptyset$ by
definition. Hence $U'\cap N''=\emptyset$. As $\Hout\subseteq U$, we get
$U'\cup N'=(U\setminus\Hout)\cup N\cup\Hout=U\cup N$. Therefore
$F_1=E_H(U',V\setminus(U\cup N))\subseteq E_H(U,V\setminus(U\cup N))=F_0$, and
since $N'\subseteq N''$, $F''=E_H(U',V\setminus(U'\cup N''))\subseteq
E_H(U',V\setminus(U'\cup N'))=F_1$. Finally $F_0\subseteq F$ by
Definition~\ref{s2:defWitness}.

(b) The vertex sets follow from \eqref{s2:eqSplit}, and the statement about
where vertices go follows from $U'\cap N''=\emptyset$. Let $e=xy\in E(H)$. If
$x,y\in U'\cup N''$ then $e\in E(G_1)$, and $e\notin F''$. Otherwise one end,
say $y$, lies outside $U'\cup N''$. If $x\in U'$ then $e\in F''$. If
$x\notin U'$ then both ends lie in $V\setminus U'$ and not both lie in $N''$,
so $e\in E(G_2)$; and $e\notin F''$ because $e$ does not meet $U'$. An edge in
$E(G_1)\cap E(G_2)$ would have both ends in $(U'\cup N'')\cap(V\setminus
U')=N''$, which \eqref{s2:eqSplit} excludes from $G_2$. So the three sets are
disjoint and cover $E(H)$.

(c) $F_0$ is determined by $(U,N)$, and every later object in (1)--(2) and in
\eqref{s2:eqSplit} is defined from $U$, $N$ and $F_0$ alone. By the Fact of
Definition~\ref{s2:defWitness}, $(U,F_0)$ is a witness with the same $N$.
\end{proof}

\begin{lemma}[Lemma SEP]\label{s2:lemSEP}
A \emph{split recursion} on a graph $H_0$ is a finite rooted binary tree whose
nodes $\nu$ carry graphs $H_\nu$, with $H_\nu=H_0$ at the root, such that at
every non-leaf node $\nu$ there are disjoint sets $U'_\nu,N''_\nu\subseteq
V(H_\nu)$ for which the two children $\nu_1,\nu_2$ of $\nu$ carry
\[
H_{\nu_1}=H_\nu[U'_\nu\cup N''_\nu],\qquad
H_{\nu_2}=H_\nu[V(H_\nu)\setminus U'_\nu]-E\bigl(H_\nu[N''_\nu]\bigr),
\]
and the \emph{deleted set} at $\nu$ is $F''_\nu\defeq
E_{H_\nu}(U'_\nu,V(H_\nu)\setminus(U'_\nu\cup N''_\nu))$. (By
\eqref{s2:eqSplit}, every split by the $\tau$-rules has this form.) We write
$|\nu|\defeq|H_\nu|$, say that a vertex $v$ \emph{lies in} $\nu$ if $v\in
V(H_\nu)$, and call the leaf nodes \emph{leaves} (distinct leaf nodes are
distinct leaves even if their vertex sets coincide). A \emph{deleted edge} is an
edge of some $F''_\nu$. Let $\Dup$ be the set of vertices lying in at least two
leaves. Then:
\begin{enumerate}
\item[(0)] at every non-leaf $\nu$, $E(H_\nu)=E(H_{\nu_1})\sqcup
E(H_{\nu_2})\sqcup F''_\nu$; a vertex of $U'_\nu$ lies only in $\nu_1$, a
vertex of $N''_\nu$ lies in both children, every other vertex of $\nu$ lies
only in $\nu_2$, and $|\nu_1|+|\nu_2|=|\nu|+|N''_\nu|$. Every vertex of a node
lies in at least one leaf below it, and the total size of the leaves is
$|H_0|+\sum_{\nu}|N''_\nu|$, the sum over the non-leaf nodes;
\item[(i)] every edge $xy$ of $H_0$ either lies in exactly one leaf, which
contains $x$ and $y$, or is deleted at exactly one node $\nu$. In the second
case $x$ and $y$ lie in $\nu$ and go to different children of $\nu$, each to
only one child, and $\nu$ is the first node on the path of nodes having $xy$
as an edge at which this happens;
\item[(ii)] if $u\notin\Dup$ and $u\in V(H_0)$, then the nodes in which $u$
lies are exactly the ancestors of the unique leaf $\Leaf_u$ containing $u$
($\Leaf_u$ included), and $u\notin N''_\nu$ at every node $\nu$ on that path;
\item[(iii)] if $u\notin\Dup$, $u,h\in V(H_0)$, $h\notin V(\Leaf_u)$, and $\nu^*$ is the
deepest ancestor of $\Leaf_u$ in which $h$ lies, then every deleted edge
$hu'$ with $u'\in V(\Leaf_u)\setminus\Dup$ is deleted at $\nu^*$.
\end{enumerate}
\deps{\ref{s2:defTauRules}}
\end{lemma}

\begin{proof}
(0) Fix a non-leaf $\nu$ and write $H=H_\nu$, $V=V(H)$, $U'=U'_\nu$,
$N''=N''_\nu$. The proof of Fact~(b) of Definition~\ref{s2:defTauRules} used
only that $U'\cap N''=\emptyset$ and the formulas \eqref{s2:eqSplit}, so it
gives $E(H)=E(H_{\nu_1})\sqcup E(H_{\nu_2})\sqcup F''_\nu$ and the statement
about where vertices go. As $U'\cap N''=\emptyset$,
$|\nu_1|+|\nu_2|=|U'|+|N''|+|V|-|U'|=|\nu|+|N''|$. Every vertex of $\nu$ lies
in a child of $\nu$, so by induction on the height of the subtree every vertex
of a node lies in some leaf below it. Summing $|\nu_1|+|\nu_2|-|\nu|=|N''_\nu|$
over all non-leaf nodes gives the total leaf size.

(i) Let $\Omega_{xy}$ be the set of nodes having $xy$ as an edge; the root is in $\Omega_{xy}$.
Every edge of a child is an edge of its parent, so $\Omega_{xy}$ is closed under taking
parents. By (0), at a non-leaf $\nu\in\Omega_{xy}$ the edge $xy$ lies in exactly one of
$E(H_{\nu_1})$, $E(H_{\nu_2})$ and $F''_\nu$; so at most one child of $\nu$ is
in $\Omega_{xy}$, and none is exactly when $xy\in F''_\nu$. Hence $\Omega_{xy}$ is a path from the
root ending at a node $\nu_\circ$. If $\nu_\circ$ is a leaf, then $xy$ is an
edge of exactly one leaf, namely $\nu_\circ$ (other leaves are not in $\Omega_{xy}$), and
$xy$ is deleted nowhere (at nodes of $\Omega_{xy}\setminus\{\nu_\circ\}$ it passes to a
child; at nodes outside $\Omega_{xy}$ it is not an edge). If $\nu_\circ$ is not a leaf,
then $xy\in F''_{\nu_\circ}$, $xy$ lies in no leaf, and it is deleted at no
other node. At $\nu_\circ$, one end lies in $U'_{\nu_\circ}$ and goes only to
the first child, and the other lies outside $U'_{\nu_\circ}\cup
N''_{\nu_\circ}$ and goes only to the second. At every earlier node of $\Omega_{xy}$ the
edge $xy$ passes to a child, which contains both $x$ and $y$.

(ii) By (0), $u$ lies in at least one leaf, hence in exactly one, $\Leaf_u$.
The nodes in which $u$ lies form a set closed under taking parents (vertex sets
shrink downwards), and at each non-leaf node of this set at least one child
contains $u$. If some node $\nu$ had both children containing $u$, then $u$
would lie in a leaf below each child, i.e.\ in two distinct leaves. So exactly
one child of every such non-leaf node contains $u$: the nodes containing $u$
form a path from the root to a leaf, which is $\Leaf_u$. A vertex of
$N''_\nu$ lies in both children of $\nu$, so $u\notin N''_\nu$ on this path.

(iii) As $h\in V(H_0)$, the root is an ancestor of $\Leaf_u$ in which $h$
lies, so $\nu^*$ exists. Let $u'\in V(\Leaf_u)\setminus\Dup$. By (ii) applied to $u'$, whose
unique leaf is $\Leaf_u$, the nodes containing $u'$ are exactly the ancestors
of $\Leaf_u$. Suppose $hu'$ is deleted at $\nu$. Then $h$ and $u'$ lie in
$\nu$, so $\nu$ is an ancestor of $\Leaf_u$ in which $h$ lies, and $\nu$ is
not below $\nu^*$. By (i), at $\nu$ the vertices $h$ and $u'$ go to different
children, each to only one. The vertex $u'$ goes to the child of $\nu$ on the
path to $\Leaf_u$, so $h$ does not lie in that child. Hence $\nu$ is the
deepest ancestor of $\Leaf_u$ containing $h$, that is, $\nu=\nu^*$.
\end{proof}

\begin{lemma}[thin cut]\label{s2:lemThinCut}
Let $\tau>0$ and consider a split recursion (Lemma~\ref{s2:lemSEP}) in which
every split is produced by the $\tau$-rules at threshold $\tau$: at each
non-leaf node $\nu$ there are $s_\nu<\tau$ and a witness at $H_\nu$ (with
parameter $s_\nu$) whose $\tau$-rules (Definition~\ref{s2:defTauRules}) give
$U'_\nu$ and $N''_\nu$. Then for every leaf $\Leaf$ and every vertex $h$,
\[
\#\bigl\{\text{deleted edges } hu:\ u\in V(\Leaf)\setminus\Dup\bigr\}\le
\lceil\tau\rceil-1,
\]
which is $\tau-1$ when $\tau$ is an integer (as in all applications), and this
number is $0$ if $h\in V(\Leaf)$. Moreover, if $u\notin\Dup$ then
every edge at $u$ is either deleted or an edge of the unique leaf containing
$u$.
\deps{\ref{s2:lemSEP}, \ref{s2:defTauRules}}
\end{lemma}

\begin{proof}
The condition $s_\nu<\tau$ only ensures that the $\tau$-rules are defined
(Definition~\ref{s2:defTauRules} asks for $\tau>s$); the argument below uses only
the structure of the $\tau$-rules.
Fix $\Leaf$ and $h$; we may assume $h$ lies in the root (otherwise no edge
of $H_0$ meets $h$) and that some deleted
edge $hu$ with $u\in V(\Leaf)\setminus\Dup$ exists. For such $u$ we have
$\Leaf_u=\Leaf$, so by Lemma~\ref{s2:lemSEP}(ii) the nodes containing $u$ are
the ancestors of $\Leaf$.

\emph{Case $h\in V(\Leaf)$.} Let $hu$ be deleted at $\nu$. Then $\nu$ is an
ancestor of $\Leaf$, and by Lemma~\ref{s2:lemSEP}(i) the vertices $h$ and $u$
go to different children of $\nu$, each to only one. But the child of $\nu$ on
the path to $\Leaf$ contains $V(\Leaf)\ni h,u$. This is a contradiction, so
the number is $0$.

\emph{Case $h\notin V(\Leaf)$.} Let $\nu^*$ be the deepest ancestor of
$\Leaf$ in which $h$ lies. By Lemma~\ref{s2:lemSEP}(iii), every deleted edge
$hu$ with $u\in V(\Leaf)\setminus\Dup$ is deleted at $\nu^*$, i.e.\ lies in
$F''=F''_{\nu^*}$. Use the notation of Definition~\ref{s2:defTauRules} at
$\nu^*$. The vertex $h$ does not lie in the child of $\nu^*$ on the path to
$\Leaf$ (the \emph{path child}), so $h\notin N''$, since vertices of $N''$ lie
in both children.

\emph{Case A: the path child is $G_1$.} Then $h\notin V(G_1)=U'\cup N''$; in
particular $h\notin U'\cup N'$ and $h\notin\Hin$, so by the definition of
$\Hin$ we have $\deg_{F_1}(h)<\tau$, i.e.\ $\deg_{F_1}(h)\le\lceil\tau\rceil-1$.
As $F''\subseteq F_1$, at most $\lceil\tau\rceil-1$ edges of $F''$ meet $h$.

\emph{Case B: the path child is $G_2$.} Then $h\notin V(G_2)=V\setminus U'$,
so $h\in U'=U\setminus\Hout$, and by the definition of $\Hout$,
$\deg_{F_0}(h)<\tau$. As $F''\subseteq F_0$, at most $\lceil\tau\rceil-1$ edges
of $F''$ meet $h$.

In both cases the count is at most $\lceil\tau\rceil-1$. The last assertion is
Lemma~\ref{s2:lemSEP}(i): an edge $uy$ that is not deleted lies in exactly one
leaf, which contains $u$ and is therefore the unique leaf $\Leaf_u$.
\end{proof}

\subsection{Overlap, and Lemma 14\texorpdfstring{$^\tau$}{-tau}}

\begin{lemma}[generic overlap bound, Lemma OV]\label{s2:lemOVgeneric}
Let $c>0$ with $3.42\,c\,\eps<1$, and put $\cOV\defeq\log_2(4/3)$. Consider a
split recursion (Lemma~\ref{s2:lemSEP}) on a root graph with $n_0\ge1$
vertices such that at every non-leaf node $\nu$, with $m\defeq|\nu|$,
\[
m\ge2,\qquad |\nu_1|\le\tfrac34m,\qquad
|N''_\nu|\le\frac{c\,\eps\,|U'_\nu|}{\log^2m}.
\]
(Recall that vertices of $U'_\nu$ go only to the first child $\nu_1$.) Let $S$
be the total size of the leaves. For $M\ge2$ let
$\Delta_{\ge M}\defeq\sum_\nu|N''_\nu|$, the sum over the non-leaf nodes $\nu$
with $|\nu|\ge M$ (the number of \emph{duplication events} at nodes of size at
least $M$; this is a number, not to be confused with the vertex set $\Dup$ of
Lemma~\ref{s2:lemSEP}), and for a vertex $v$ let $\dup_{\ge M}(v)$ be the number of
non-leaf nodes $\nu$ with $|\nu|\ge M$ and $v\in N''_\nu$ (so
$\sum_v\dup_{\ge M}(v)=\Delta_{\ge M}$). Then:
\begin{enumerate}
\item[(a)] $\displaystyle\Delta_{\ge M}\le
c\,\eps\Bigl(\frac1{\log^2M}+\frac1{\cOV\log M}\Bigr)S\le
\frac{3.42\,c\,\eps\,S}{\log M}$;
\item[(b)] for every vertex $v$, the number of leaves $\Leaf$ with
$|\Leaf|\ge M$ containing $v$ is at most $1+\dup_{\ge M}(v)$; consequently
$\sum_{\Leaf:|\Leaf|\ge M}|\Leaf|-\bigl|\bigcup_{\Leaf:|\Leaf|\ge
M}V(\Leaf)\bigr|\le\Delta_{\ge M}$;
\item[(c)] $S\le n_0/(1-3.42\,c\,\eps)$.
\end{enumerate}
\emph{Instance $c=1$ (the $s=0$ recursion).} Call a split recursion an
\emph{$s=0$ recursion} if at every non-leaf node $\nu$ there is a witness
$(U_\nu,F_\nu)$ at $H_\nu$ with parameter $s=0$ (Definition~\ref{s2:defWitness})
and $U'_\nu=U_\nu$, $N''_\nu=\Nbr_{H_\nu}(U_\nu)$. This is exactly the
recursion in the proof of \cite[Lemma~14]{BM} with $s=0$
(Cited result~\ref{s1:citLem14}): there $F_\nu=\emptyset$, the children are
$H_\nu[U_\nu\cup N_\nu]$ and $H_\nu\setminus U_\nu-E(H_\nu[U_\nu\cup
N_\nu])=H_\nu[V(H_\nu)\setminus U_\nu]-E(H_\nu[N_\nu])$, and the deleted set
$E_{H_\nu}(U_\nu,V(H_\nu)\setminus(U_\nu\cup N_\nu))$ is empty. (It is also
the split by the $\tau$-rules at $s=0$ for any $\tau>0$: then $F_0=\emptyset$
and $\Hout=\Hin=\emptyset$.) Every $s=0$ recursion satisfies the hypotheses
with $c=1$; hence its total leaf size is at most $n_0/(1-3.42\eps)\le1.12\,n_0$.
The lemma is applied with $c=1.6$ to the recursion of Lemma~14$^\tau$ below,
and with $c=1.6$ to the two-level recursion of a round of $\HBtp$.
\deps{\ref{s2:lemSEP}, \ref{s2:defWitness}, \ref{s1:citLem14}}
\end{lemma}

\begin{proof}
(a) If no non-leaf node has size at least $M$ there is nothing to prove. For
every non-leaf node $\nu$ with $|\nu|\ge M$ and every $u\in U'_\nu$, give the
pair $(\nu,u)$ the charge $c\eps/\log^2|\nu|$. The total charge at $\nu$ is
$c\eps|U'_\nu|/\log^2|\nu|\ge|N''_\nu|$, so $\Delta_{\ge M}$ is at most the
total charge. A \emph{leaf copy} is a pair $(\Leaf,v)$ with $\Leaf$ a leaf and
$v\in V(\Leaf)$; there are exactly $S$ leaf copies. For each charged pair
$(\nu,u)$, the vertex $u$ lies in $\nu_1$, hence (Lemma~\ref{s2:lemSEP}(0)) in
some leaf $\Leaf$ below $\nu_1$; fix one and map $(\nu,u)$ to the leaf copy
$(\Leaf,u)$. The pair $(\nu,u)$ is recovered from $u$ and $\nu$, and $\nu$ is
an ancestor of $\Leaf$; so each leaf copy $(\Leaf,u)$ receives the charges of
pairs $(\nu,u)$ for a set of ancestors $\nu$ of $\Leaf$ with $|\nu|\ge M$,
$u\in U'_\nu$ and $\Leaf$ below $\nu_1$. List these ancestors as
$\nu^{(1)},\nu^{(2)},\dots,\nu^{(k)}$ from $\Leaf$ upwards. For $i<k$, the
node $\nu^{(i)}$ lies strictly below $\nu^{(i+1)}$ on the path from
$\nu^{(i+1)}$ to $\Leaf$, which passes through the first child of
$\nu^{(i+1)}$; vertex sets shrink downwards, so
$|\nu^{(i)}|\le|\nu^{(i+1)}_1|\le\frac34|\nu^{(i+1)}|$. Hence
$|\nu^{(i)}|\ge(4/3)^{i-1}|\nu^{(1)}|\ge(4/3)^{i-1}M$, i.e.\
$\log|\nu^{(i)}|\ge\log M+(i-1)\cOV$. The charge received by $(\Leaf,u)$ is
therefore at most
\[
c\eps\sum_{j\ge0}\frac1{(\log M+j\cOV)^2}\le
c\eps\Bigl(\frac1{\log^2M}+\int_0^\infty\frac{dx}{(\log M+x\cOV)^2}\Bigr)
=c\eps\Bigl(\frac1{\log^2M}+\frac1{\cOV\log M}\Bigr),
\]
because the summand is decreasing in $j$. Summing over the $S$ leaf copies
gives the first bound. For $M\ge2$ we have $\log M\ge1$, so
$1/\log^2M\le1/\log M$, and $1+1/\cOV=3.4095\ldots\le3.42$.

(b) Let $\mathcal{C}_v$ be the set of nodes of size at least $M$ in which $v$ lies. It
is closed under taking parents (vertex sets, hence sizes, shrink downwards), so
if nonempty it is a rooted subtree containing the root. Every leaf of the
recursion of size at least $M$ that contains $v$ belongs to $\mathcal{C}_v$ and has no
children, so it is a leaf of $\mathcal{C}_v$. In a finite rooted tree the number of
leaves is $1+\sum_\nu(\kappa_\nu-1)$, the sum over the nodes with
$\kappa_\nu\ge1$ children; here $\kappa_\nu\le2$, and $\kappa_\nu=2$ means that
$v$ lies in both children of $\nu$, i.e.\ $v\in N''_\nu$
(Lemma~\ref{s2:lemSEP}(0)), with $|\nu|\ge M$. So the number of such leaves is
at most $1+\dup_{\ge M}(v)$. Summing $(\#\{\Leaf\ni v:|\Leaf|\ge M\}-1)^+\le
\dup_{\ge M}(v)$ over $v$ gives the consequence.

(c) All non-leaf nodes have size at least $2$, so by
Lemma~\ref{s2:lemSEP}(0), $S=n_0+\sum_\nu|N''_\nu|=n_0+\Delta_{\ge2}\le
n_0+3.42c\eps S$ by (a) with $M=2$. Rearranging gives (c).

\emph{Instance $c=1$.} By Definition~\ref{s2:defWitness} with $s=0$ we have
$F_\nu=\emptyset$, so $N_\nu=\Nbr_{H_\nu}(U_\nu)$, and every neighbour of
$U_\nu$ outside $U_\nu$ lies in $N_\nu$; this gives the three formulas stated.
The node has $m\ge2$ vertices (Definition~\ref{s2:defWitness}), and
$|N''_\nu|=|N_\nu|<\eps|U_\nu|/\log^2m$, which is the hypothesis with $c=1$;
and $|\nu_1|=|U_\nu|+|N_\nu|<(1+\eps)\frac23m<\frac34m$, using $\log m\ge1$.
Finally $1/(1-3.42/32)=1.1196\ldots\le1.12$.
\end{proof}

\begin{lemma}[Lemma 14$^\tau$]\label{s2:lem14tau}
Let $H_0$ be a graph on $n_0$ vertices, $s\ge1$ an integer, $\eps=2^{-5}$, and
$\tau\ge128\,s\log^2n_0$. The \emph{$\tau$-run} of $H_0$ (with parameters
$s,\tau$) is the recursion that, at a node $H$, stops if $H$ is an
$(\eps,s)$-expander (the node is then a leaf), and otherwise picks any witness
$(U,F)$ at $H$ with parameter $s$ (Definition~\ref{s2:defWitness}) and splits
$H$ by the $\tau$-rules at threshold $\tau$ (Definition~\ref{s2:defTauRules};
note $\tau\ge128s>s$ whenever $n_0\ge2$). Let $\Dup$ be the set of vertices
lying in at least two leaves. Then, for every choice of witnesses:
\begin{enumerate}
\item[(a)] at every split, at a node of size $m$,
\[
|\Hout|+|\Hin|\le\frac{|F_0|}{\tau}\le\frac{s|U|}{\tau}\le\frac{|U|}{128\log^2m},
\qquad |N''|<\frac{1.25\,\eps|U|}{\log^2m}<\frac{1.5\,\eps|U|}{\log^2m},
\]
$|U'|\ge\frac{127}{128}|U|>0$, $U\subseteq U'\cup N''$, $n_1\defeq|U'\cup
N''|\le0.698\,m<\frac34m$ and $n_2\defeq m-|U'|<m$. The recursion terminates;
it is a split recursion; at most $4sn_0\log n_0$ edges are deleted; the
remaining edges are partitioned into the leaves; every leaf is an
$(\eps,s)$-expander; and the total size of the leaves is at most
$2n_0-2n_0/(2+\log n_0)$;
\item[(b)] (OV at $1.6\eps$) at every split, $|N''|\le1.6\,\eps|U'|/\log^2m$,
the vertices of $U'$ go only to $G_1$, and $|G_1|\le\frac34m$. Hence, by
Lemma~\ref{s2:lemOVgeneric} with $c=1.6$, the total leaf size $S$ satisfies
$S\le n_0/(1-5.5\eps)\le1.21\,n_0$, and for every $M\ge2$ the duplication at
nodes of size at least $M$ is $\Delta_{\ge M}\le5.46\,\eps S/\log M$;
\item[(c)] for every leaf $\Leaf$ and every vertex $h$,
$\#\{\text{deleted edges }hu:u\in V(\Leaf)\setminus\Dup\}\le\lceil\tau\rceil-1$
($=\tau-1$ for integer $\tau$), and the number is $0$ if $h\in V(\Leaf)$;
\item[(d)] if $u\notin\Dup$, every edge at $u$ is deleted or lies in the unique
leaf containing $u$.
\end{enumerate}
\deps{\ref{s1:citLem14}, \ref{s1:citDef11}, \ref{s2:defTauRules}, \ref{s2:lemThinCut}, \ref{s2:lemOVgeneric}, \ref{s2:lemSEP}, \ref{s2:defWitness}}
\end{lemma}

\begin{proof}
If $n_0=1$ the root is an $(\eps,s)$-expander and everything is trivial; let
$n_0\ge2$. Every node of the recursion is a subgraph of $H_0$ on $m\le n_0$
vertices, so $\tau\ge128\,s\log^2m$ at every node. The argument follows the
proof of \cite[Lemma~14]{BM} (Cited result~\ref{s1:citLem14}) with the split
of Definition~\ref{s2:defTauRules} in place of the split of \cite{BM}.

\emph{Counting at one split.} Fix a non-leaf node $H$ on $m\ge2$ vertices,
with witness $(U,F)$, and use the notation of
Definition~\ref{s2:defTauRules}; put $V\defeq V(H)$. Every edge of $F_0$ has
exactly one end in $U$, so $\sum_{v\in U}\deg_{F_0}(v)=|F_0|$. Let
$a\defeq\sum_{v\in\Hout}\deg_{F_0}(v)$; then $\tau|\Hout|\le a$. By Fact~(a) of
Definition~\ref{s2:defTauRules},
$F_1=E_H(U',V\setminus(U\cup N))$ is $F_0$ minus the edges at $\Hout$, so
$|F_1|=|F_0|-a$. Every edge of $F_1$ has exactly one end outside $U'\cup
N'=U\cup N$, so $\tau|\Hin|\le\sum_{x\notin U\cup N}\deg_{F_1}(x)=|F_1|$.
Adding, $\tau(|\Hout|+|\Hin|)\le|F_0|\le|F|\le s|U|$. With
$\tau\ge128\,s\log^2m$ this gives $|\Hout|+|\Hin|\le|U|/(128\log^2m)$, and
$1/128=\eps/4$. Therefore
\[
|N''|=|N|+|\Hout|+|\Hin|<\frac{\eps|U|}{\log^2m}+\frac{\eps|U|}{4\log^2m}
=\frac{1.25\,\eps|U|}{\log^2m}.
\]
Since $|\Hout|\le|U|/128$, $|U'|\ge\frac{127}{128}|U|>0$. As
$\Hout\subseteq N'\subseteq N''$, $U=U'\cup\Hout\subseteq U'\cup N''$, so
$|U|\le n_1$. Also $n_1=|U'|+|N''|\le|U|+|N''|<(1+1.5\eps)|U|\le
\frac23(1+\frac{3}{64})m\le0.698\,m$, using $\log m\ge1$; and $n_2=|V\setminus
U'|=m-|U'|<m$. By Fact~(b) of Definition~\ref{s2:defTauRules},
$n_1+n_2=m+|N''|$. Finally $\log n_1\le\log m+\log0.698<\log m-\frac25$, since
$\log_20.698=-0.518\ldots$.

\emph{Termination.} Both children of every split have fewer vertices than
their parent, so every root-to-node path has at most $n_0$ nodes, and the
binary tree is finite. By \eqref{s2:eqSplit} the recursion is a split
recursion. Its leaves are $(\eps,s)$-expanders by the stopping rule, and by
Lemma~\ref{s2:lemSEP}(i) every edge of $H_0$ is either deleted (exactly once)
or lies in exactly one leaf.

\emph{The induction.} We show by induction on the height of the subtree below
a node $\nu$, with $m=|\nu|$, that the subtree deletes at most $4sm\log m$ edges
and its leaves have total size at most $2m-2m/(2+\log m)$. For a leaf this
holds, since nothing is deleted and $m\le2m-2m/(2+\log m)$ (as
$2+\log m\ge2$). Let $\nu$ be a non-leaf node, with $n_1,n_2$ as above
($1\le n_1,n_2<m$), and let $E^{\mathrm{del}}_i$ be the set of edges deleted in the subtree
below the $i$-th child; by induction $|E^{\mathrm{del}}_i|\le4sn_i\log n_i$. The edges
deleted in the subtree below $\nu$ are $F''\cup E^{\mathrm{del}}_1\cup E^{\mathrm{del}}_2$, and $|F''|\le
|F_0|\le s|U|\le sn_1$. As $|U|\le n_1$, we have $|N''|<1.5\eps n_1/\log^2m$.
Hence, as in \cite[inequality~(8)]{BM},
\begin{align*}
\frac{|F''|+|E^{\mathrm{del}}_1|+|E^{\mathrm{del}}_2|}{s}
&\le n_1+4n_1\log n_1+4n_2\log n_2
\le n_1+4n_1\bigl(\log m-\tfrac25\bigr)+4n_2\log m\\
&=4(n_1+n_2)\log m-\tfrac35n_1=4(m+|N''|)\log m-\tfrac35n_1\\
&<4m\log m+n_1\Bigl(\frac{6\eps}{\log m}-\frac35\Bigr)\le4m\log m,
\end{align*}
since $6\eps/\log m\le6/32<3/5$. For the leaf sizes put $L\defeq\log m\ge1$.
As $n_2<m$, $2n_2/(2+\log n_2)\ge2n_2/(2+L)$. As $\log n_1<L-\frac25$,
\[
\frac{2n_1}{2+\log n_1}>\frac{2n_1}{\frac85+L}=\frac{2n_1}{2+L}
+\frac{\frac45n_1}{(\frac85+L)(2+L)}\ge\frac{2n_1}{2+L}+\frac{n_1}{10L^2},
\]
where the last step is $8L^2\ge(\frac85+L)(2+L)$, i.e.\
$7L^2-3.6L-3.2\ge0$, which holds for $L\ge1$ (the left side is $0.2$ at
$L=1$ and increasing for $L\ge1$). By induction the total leaf size below
$\nu$ is therefore
\begin{align*}
\sum_{i=1}^2\Bigl(2n_i-\frac{2n_i}{2+\log n_i}\Bigr)
&<(n_1+n_2)\Bigl(2-\frac2{2+L}\Bigr)-\frac{n_1}{10L^2}
=(m+|N''|)\Bigl(2-\frac2{2+L}\Bigr)-\frac{n_1}{10L^2}\\
&\le2m-\frac{2m}{2+L}+2|N''|-\frac{n_1}{10L^2}\\
&<2m-\frac{2m}{2+L}+\Bigl(3\eps-\frac1{10}\Bigr)\frac{n_1}{L^2}
\le2m-\frac{2m}{2+L},
\end{align*}
as in \cite[inequality~(9)]{BM}, because $3\eps=3/32<1/10$. At the root this gives (a).

(b) By (a), $|N''|<1.5\eps|U|/\log^2m\le1.5\cdot\frac{128}{127}\eps|U'|/\log^2m
<1.52\,\eps|U'|/\log^2m$. (The first bound of (a) even gives
$|N''|<1.25\cdot\frac{128}{127}\eps|U'|/\log^2m<1.26\,\eps|U'|/\log^2m$; the
looser constant $1.6$ suffices.) Vertices of $U'$ go only to $G_1$ and $|G_1|=n_1\le
\frac34m$ (Fact~(b) of Definition~\ref{s2:defTauRules} and (a)); every
non-leaf node has $m\ge2$. As $3.42\cdot1.6\,\eps=0.171<1$,
Lemma~\ref{s2:lemOVgeneric} applies with $c=1.6$: by its part (c),
$S\le n_0/(1-5.472\eps)\le n_0/(1-5.5\eps)=1.2075\ldots\,n_0\le1.21\,n_0$; by its
part (a), $\Delta_{\ge M}\le1.6\eps(1+1/\cOV)S/\log M\le5.46\,\eps S/\log M$,
since $1.6\,(1+1/\cOV)=5.455\ldots$.

(c), (d) Every split is produced by the $\tau$-rules at threshold $\tau$ from a
witness with parameter $s<\tau$, so Lemma~\ref{s2:lemThinCut} applies.
\end{proof}

\subsection{The hierarchy}

\begin{definition}[the hierarchy $\HBtp$]\label{s2:defHBtp}
Let $G$ be a graph on $n$ vertices, and let $\eps=2^{-5}$, $\sigma=100$,
$\Cp=103$, $\Aexp=105$ and $\Dstar$ be as in Definition~\ref{s1:defConstants}.
Put $G_1\defeq G$. All graphs $G_l$, $G'_l$ below have vertex set $V(G)$. For
$l=1,2,\dots$:
\begin{enumerate}
\item[(R0)]\namedlabel{s2:R0}{(R0)} $d_l\defeq2|E(G_l)|/n$. If $d_l<\Dstar$,
stop, and put $E_0\defeq E(G_l)$ and $R\defeq l-1$. Otherwise put
$\lambda_l\defeq\log d_l$.
\item[(R1)]\namedlabel{s2:R1}{(R1)} $\tHB_l\defeq\log^2d_l$. As long as the
current graph contains a cycle of length at least $\tHB_ld_l$, delete one (the
lexicographically first, for a fixed order of $V(G)$). The deleted cycles form
$\Cyc_l$; the remaining graph is $G'_l$. Thus $G'_l$ has no cycle of length at
least $\tHB_ld_l$.
\item[(R2)]\namedlabel{s2:R2}{(R2)}
$M_l\defeq\bigl\lceil\max\bigl(2^{40},\,2^{16}\,\tHB_ld_l\log^4(\tHB_ld_l)\bigr)\bigr\rceil\in\mathbb N$,\quad
$\Lambda_l\defeq\log M_l$,\quad $s_l\defeq\lceil\Lambda_l^{\sigma}\rceil$,\quad
$P_l\defeq\lceil\lambda_l^{\Cp}\rceil$,\quad
$\tau_l\defeq\lceil128\,s_l\log^2M_l\rceil$.
The ceiling makes $M_l$ a natural number, as its later uses require (it counts
classes and colours, for instance $\KJS_l=M_l^2$ in Definition~\ref{s3:defCOL}
and the palettes $[3M_l]$, $[4M_l]$ of Section~\ref{s7:sec}); $\Lambda_l$ is the
logarithm of this integer. Besides integrality, only two facts about the ceiling are used:
$M_l\ge\max\bigl(2^{40},\,2^{16}\,\tHB_ld_l\log^4(\tHB_ld_l)\bigr)$, and
$M_l\le\max\bigl(2^{40},\,2^{16}\,\tHB_ld_l\log^4(\tHB_ld_l)\bigr)+1$.
\item[(R3)]\namedlabel{s2:R3}{(R3)} Run an $s=0$ recursion
(Lemma~\ref{s2:lemOVgeneric}) on $G'_l$ that stops exactly at
$(\eps,0)$-expanders: a node is a leaf iff its graph is an
$(\eps,0)$-expander, and otherwise any witness with parameter $0$ is used. This
is the recursion of the proof of \cite[Lemma~14]{BM} with $s=0$ (Cited
result~\ref{s1:citLem14}). Its leaves are the \emph{$s=0$ pieces} $\piece$ of
round $l$. Inside every piece with $|\piece|\ge P_l$, run the $\tau$-run of
$\piece$ with $s=s_l$ and $\tau=\tau_l$ (Lemma~\ref{s2:lem14tau}), with any
witnesses. The \emph{two-level recursion} of round $l$ is the split recursion
obtained from the $s=0$ recursion by attaching, at every piece $\piece$ with
$|\piece|\ge P_l$, the tree of its $\tau$-run. Its leaves are the pieces of
size less than $P_l$ and the leaves of the $\tau$-runs. The \emph{round-$l$
pre-parts} are the leaves of the $\tau$-runs with at least $P_l$ vertices;
equivalently, the leaves of the two-level recursion with at least $P_l$
vertices. For a pre-part named $Z$, $Z^0$ denotes its vertex set and $X^0_Z$
its leaf graph, a graph on $Z^0$. Finally
$\Dups_l\defeq$ the set of vertices lying in at least two leaves (of any size,
singletons included) of the two-level recursion of round $l$.
\item[(R4)]\namedlabel{s2:R4}{(R4)} $D_l\defeq$ the set of vertices lying in at
least two round-$l$ pre-parts. Fix an order of the round-$l$ pre-parts (the
\emph{home order}). For a vertex $v$ lying in some pre-part,
$\home(v)=\home_l(v)\defeq$ the first pre-part containing $v$. The
\emph{guests} of a pre-part $Z$ are $S_Z\defeq\{v\in Z^0\cap D_l:\home(v)\ne
Z\}$. The pre-part $Z$ is \emph{light} iff
\begin{enumerate}
\item[(L1)] $|S_Z|\le|Z^0|/2$;
\item[(L2)] every $u\in Z^0$ has at most $s_l/2$ neighbours in $S_Z$ in the
graph $G'_l[Z^0]$;
\item[(GC)]\namedlabel{s2:ruleGC}{(GC)} no $x\in S_Z$ has at least
\[
\thGC_l(Z^0)\defeq\bigl\lceil|Z^0|\,\lambda_l^{-1/2}\bigr\rceil
\]
neighbours in $Z^0\setminus S_Z$ in the graph $X^0_Z$ (the \emph{guest cap}).
\end{enumerate}
Otherwise $Z$ is \emph{standalone}. A pre-part satisfying (L1) and (L2) but
failing (GC) is a \emph{GC-part}; it is standalone. A light pre-part $Z$ gives
the \emph{light part} with vertex set $Z^0\setminus S_Z$ and graph
$X_Z\defeq X^0_Z-S_Z$. $\Std_l$ is the set of round-$l$ standalone pre-parts;
it contains the GC-parts and the pre-parts failing (L1) or (L2), and it is a
deterministic function of the run.
\item[(R5)]\namedlabel{s2:R5}{(R5)} The edges of $G'_l$ are assigned in this
order.
(1) For every light pre-part $Z$ the edges of $X_Z$ are assigned to the light
part $Z$; for every standalone pre-part $Z$ the edges of $X^0_Z$ are assigned
to $Z$. (The graphs $X^0_Z$ of distinct pre-parts are edge-disjoint, being
distinct leaves of the two-level recursion, Lemma~\ref{s2:lemSEP}(i), and
$X_Z\subseteq X^0_Z$; so no edge is assigned twice in this step.)
(2) Every edge not yet assigned whose ends both lie in some light part is
assigned to the first such light part (in the home order).
(3) Every edge not yet assigned whose ends both lie in $Z^0$ for some
standalone pre-part $Z$ is assigned to the first such pre-part.
(4) Every other edge of $G'_l$ passes down: $E(G_{l+1})$ is the set of these
edges.
$E_l(Z)$ denotes the set of edges assigned at round $l$ to $Z$ (a light part or
a standalone pre-part).
\end{enumerate}
\emph{Naming.} Pre-parts are leaves, so distinct pre-parts are counted
separately even if their vertex sets coincide. If a pre-part $Z$ is light, the
letter $Z$ also denotes the vertex set $Z^0\setminus S_Z$ of its light part.
A \emph{round-$l$ part} is a light part or a standalone pre-part of round $l$.
We write $\home_l$, $S_Z$ and so on with the round when it is not clear from
the context.

A \emph{valid $\HBtp$ run} is an execution of this procedure with any
witnesses in both levels of (R3), any vertex order in (R1), and any home
orders in (R4). At $s=0$ the $\tau$-rules are idle
(Lemma~\ref{s2:lemOVgeneric}), so the first level of (R3) is the recursion
of \cite{BM} unchanged; the $\tau$-rules act only inside the pieces.
\deps{\ref{s1:citLem14}, \ref{s2:lem14tau}, \ref{s1:defConstants}, \ref{s2:lemOVgeneric}, \ref{s2:lemSEP}}
\fixes{\RTa{M10}, \RTb{I8}}
\end{definition}

\begin{definition}[ancestors, typing, multiplicities]\label{s2:defAncestors}
Fix a valid $\HBtp$ run. An \emph{ancestor of round $r$} is either
\begin{itemize}
\item a light part $Y$ of round $r$; then $V(Y)\defeq Y$ (the vertex set
$Y^0\setminus S_Y$), $H_Y\defeq X_Y$, $\eps_Y\defeq2^{-6}$ and
$s_Y\defeq s_r/2$; or
\item a standalone pre-part $Y\in\Std_r$, GC-parts included; then
$V(Y)\defeq Y^0$, $H_Y\defeq X^0_Y$, $\eps_Y\defeq2^{-5}$ and $s_Y\defeq s_r$.
\end{itemize}
We write $r(Y)$ for the round of $Y$ and $L_Y\defeq\log|V(Y)|$. Every ancestor
corresponds to a distinct pre-part $Y^0\supseteq V(Y)$ of its round. For
$l\ge3$ and a vertex $x$,
\[
\anc_l(x)\defeq\{Y:\ Y\text{ an ancestor of a round }r\le l-2\text{ with }x\in
V(Y)\},
\]
and $\anc_l(x)\defeq\emptyset$ for $l\le2$. For $Z\in\Std_l$:
\begin{itemize}
\item the \emph{hubs} of $Z$ are $A_Z\defeq Z^0\cap D_l$;
\item the \emph{ports} of $Z$ are $U_Z\defeq Z^0\setminus D_l$;
\item the \emph{fresh ports} are $F_Z\defeq\{x\in U_Z:\anc_l(x)=\emptyset\}$
(so all ports are fresh for $l\le2$);
\item the \emph{classed ports} are $Q_Z\defeq U_Z\setminus F_Z$.
\end{itemize}
For a round $r$ and a vertex $w$: $\mu_r(w)$ is the number of round-$r$
pre-parts containing $w$ (so $D_r=\{w:\mu_r(w)\ge2\}$);
$\dup_r\defeq\sum_w(\mu_r(w)-1)^+$; and $\mult_r(w)$ is the number of ancestors
$Y$ of round $r$ with $w\in V(Y)$. Finally $j_0(x)$ is the first round in which
$x$ lies in a pre-part, and $j_1(x)$ the first round in which $x$ lies in a
light part ($\infty$ if there is none).
\deps{\ref{s2:defHBtp}}
\end{definition}

\subsection{Basic properties of valid runs}

\begin{lemma}[Lemma-25 size cap]\label{s2:lemCap}
\begin{enumerate}
\item[(i)] If $m\ge2^{40}$, $T\ge2^{117}$ and $m<18432\,T\log^4m$, then
$m\le2^{16}T\log^4T$.
\item[(ii)] In every round $l\le R$ of a valid $\HBtp$ run, every $s=0$ piece
$\piece$ satisfies $|\piece|\le M_l$. Hence every round-$l$ pre-part has
$|Z^0|\le M_l$; for every round-$l$ part $Z$ (light or standalone) every vertex
is incident with at most $M_l-1$ edges of $E_l(Z)$; and
$\tau_l\ge128\,s_l\log^2|\piece|$, so Lemma~\ref{s2:lem14tau} applies to the
$\tau$-run of every piece with $|\piece|\ge P_l$.
\end{enumerate}
\emph{Remark (not used).} Part (i) is false for small $T$: for $T=2^{20}$ and
$m=2^{55}$ we have $m\ge2^{40}$ and $18432\,T\log^4m=2^{57.29\ldots}>m$, but
$2^{16}T\log^4T=2^{53.28\ldots}<m$. The uniform form $m\le\max\bigl(2^{40},
2^{16}T(\log T+40)^4\bigr)$ holds for all $T\ge1$.
\deps{\ref{s1:citLem25}, \ref{s1:condG2}, \ref{s2:defHBtp}, \ref{s2:lem14tau}}
\end{lemma}

\begin{proof}
(i) Put $\chi(y)\defeq y/\log^4y$; since $\frac{d}{dy}\ln \chi(y)=\frac1y\bigl(1-
\frac4{\ln y}\bigr)$, $\chi$ is increasing for $y>e^4$. The hypothesis says
$\chi(m)<18432\,T$. Put $x\defeq\log T\ge117$ and $B\defeq2^{16}T\log^4T=2^{16}Tx^4$,
so $\log B=16+x+4\log x$. We have $(2^{16}/18432)^{1/4}=(32/9)^{1/4}=1.37317\ldots$,
and $0.37317\,x\ge16+4\log x$ for $x\ge117$: at $x=117$ the two sides are
$43.66\ldots$ and $43.48\ldots$, and for $x\ge117$ the derivative of the left
side ($0.373\ldots$) exceeds that of the right side ($4/(x\ln2)<0.05$). Hence
$(32/9)^{1/4}x\ge16+x+4\log x=\log B$, i.e.\ $2^{16}x^4\ge18432\,(\log B)^4$,
i.e.\ $\chi(B)\ge18432\,T$. If $m>B$, then, as $B>e^4$, $\chi(m)>\chi(B)\ge18432\,T$,
contradicting the hypothesis. So $m\le B$.

\emph{Remark.} The numbers in the counterexample are direct computations. For
the uniform form put $B'\defeq2^{16}T(x+40)^4$ with $x=\log T\ge0$; then $\log
B'=16+x+4\log(x+40)$ and $\chi(B')\ge18432\,T$ is equivalent to
$(32/9)^{1/4}(x+40)\ge16+x+4\log(x+40)$, which follows from
$0.37317\,x+38.9\ge4\log(x+40)$; this holds at $x=0$ ($38.9\ge21.3$) and the
left side grows faster ($0.373$ against at most $4/(40\ln2)<0.15$). As
$B'>e^4$, $m>\max(2^{40},B')$ would give $\chi(m)>\chi(B')\ge18432\,T$.

(ii) Let $l\le R$ and $T\defeq\tHB_ld_l=d_l\log^2d_l$. By \Gc{2}(a)
(Condition~\ref{s1:condG2}), $d_l\ge\Dstar\ge2^{117}$, so $\log^2d_l\ge1$ and
$T\ge2^{117}$; and $M_l\ge\max(2^{40},2^{16}T\log^4T)$ by (R2). Let $\piece$ be an
$s=0$ piece and $m\defeq|\piece|$. If $m<2^{40}$ then $m\le M_l$. Otherwise the
graph of $\piece$ is an $(\eps,0)$-expander (the stopping rule in (R3)) with
$\eps=2^{-5}$ on $m\ge2^{40}=2^{30}/\eps^2$ vertices, so by Cited
result~\ref{s1:citLem25} it contains a cycle of length at least
$\eps^2m/(18\log^4m)=m/(18432\log^4m)$. This cycle lies in $G'_l$, which has no
cycle of length at least $T$ by (R1). Hence $m<18432\,T\log^4m$, and (i) gives
$m\le2^{16}T\log^4T\le M_l$. Every pre-part is a leaf of the $\tau$-run of some
piece, so $Z^0\subseteq V(\piece)$ and $|Z^0|\le M_l$. By (R5) every edge of
$E_l(Z)$ has both ends in $Z^0$ (in the light part $Z\subseteq Z^0$ if $Z$ is
light), so every vertex meets at most $|Z^0|-1\le M_l-1$ of them. Finally
$\tau_l\ge128\,s_l\log^2M_l\ge128\,s_l\log^2|\piece|$.
\end{proof}

\begin{lemma}[Lemma HS]\label{s2:lemHS}
Let $\varepsilon'>0$ and $s\ge0$, let $X$ be an $(\varepsilon',s)$-expander on
a vertex set of size $z\ge11$, let $W\subseteq V(X)$ with $|W|\le z/2$, and
suppose every $u\in V(X)\setminus W$ has at most $s_W\le s$ neighbours in $W$ in
$X$. Then $X-W$ is an $\bigl(\varepsilon'(\log z'/\log z)^2,\,s-s_W\bigr)$-expander
on $z'\defeq z-|W|$ vertices, and $\varepsilon'(\log z'/\log z)^2\ge
\varepsilon'/2$. In particular, if $X$ is a $(2^{-5},s)$-expander and
$s_W=s/2$, then $X-W$ is a spanning $(2^{-6},s/2)$-expander of $V(X)\setminus
W$.
\deps{\ref{s1:citDef11}}
\end{lemma}

\begin{proof}
Note $z'\ge z/2\ge5.5$. Let $U\subseteq V(X)\setminus W$ with $1\le|U|\le2z'/3$
and $F'\subseteq E(X-W)$ with $|F'|\le(s-s_W)|U|$. Put $F\defeq F'\cup
E_X(U,W)$. Then $|E_X(U,W)|\le s_W|U|$, so $|F|\le s|U|$, and $|U|\le2z'/3\le
2z/3$. By Cited result~\ref{s1:citDef11} applied to $X$,
$|\Nbr_{X-F}(U)|\ge\varepsilon'|U|/\log^2z$. No vertex of $W$ lies in
$\Nbr_{X-F}(U)$, because all edges between $U$ and $W$ are in $F$. For
$v\notin U\cup W$ and $u\in U$, the edge $uv$ (if present) lies in $X-W$ and not
in $E_X(U,W)$, so $uv\notin F$ iff $uv\notin F'$. Hence
$\Nbr_{X-F}(U)=\Nbr_{(X-W)-F'}(U)$, and
\[
|\Nbr_{(X-W)-F'}(U)|\ge\frac{\varepsilon'|U|}{\log^2z}
=\varepsilon'\Bigl(\frac{\log z'}{\log z}\Bigr)^2\frac{|U|}{\log^2z'}.
\]
This is the expansion claimed. Next, $(\log z'/\log z)^2\ge\frac12$ iff $z'\ge
z^{1/\sqrt2}$. Since $z'\ge z/2$, it suffices that $z^{1-1/\sqrt2}\ge2$, i.e.\
$\log z\ge1/(1-1/\sqrt2)=3.414\ldots$, i.e.\ $z\ge10.67$; this holds as $z\ge11$.
The last sentence follows because an $(\varepsilon_a,s)$-expander is an
$(\varepsilon_b,s)$-expander whenever $\varepsilon_b\le\varepsilon_a$.
\end{proof}

\begin{proposition}[structure of a valid run]\label{s2:propStructure}
For every valid $\HBtp$ run on a graph $G$ with $n\ge\Nzero$ and $d_1\ge\Dstar$:
\begin{enumerate}
\item[(i)] each $X^0_Z$ is a $(2^{-5},s_l)$-expander on $Z^0$; each light $X_Z$
is a $(2^{-6},s_l/2)$-expander on the light part $Z$; every $H_Y$ has minimum
degree greater than $s_Y$; and $s_r\ge\lambda_r^{100}$ for every $r\le R$;
\item[(ii)] the conclusions of Lemma~\ref{s2:lemCap}(ii) hold;
\item[(iii)] $E(G_{l+1})\subseteq E(G'_l)\subseteq E(G_l)$ for every $l\le R$,
so $d_{l+1}\le d_l$; and
\[
E(G)=\bigsqcup_{l\le R}E(\Cyc_l)\ \sqcup\ E_0\ \sqcup
\bigsqcup_{Y\text{ light}}E_{r(Y)}(Y)\ \sqcup\bigsqcup_{l\le R}\
\bigsqcup_{Z\in\Std_l}E_l(Z)
\]
(disjoint unions); $E(H_Y)\subseteq E_{r(Y)}(Y)$ for every ancestor $Y$; every
edge of $E_{r(Y)}(Y)$ has both ends in $V(Y)$; the graphs $H_Y$ of all
ancestors $Y$ of all rounds are pairwise edge-disjoint, and the graphs $X^0_Z$
of the pre-parts of one round are pairwise edge-disjoint; $\sum_{l\le
R}|\Cyc_l|\le n$ (the number of long cycles); and $|E_0|<\Dstar n/2$;
\item[(iv)] light parts of one round are pairwise vertex-disjoint, and
$|Y|\ge|Y^0|/2\ge P_r/2$ for every light part $Y$ of round $r$; every vertex
outside $D_l$ lies in at most one round-$l$ pre-part; the port sets $U_Z$ of
distinct $Z\in\Std_l$ are pairwise disjoint; and each vertex lies in at most one
light part per round, namely in the light part of $\home_l(v)$ if that
pre-part is light.
\end{enumerate}
\deps{\ref{s2:defHBtp}, \ref{s2:defAncestors}, \ref{s2:lem14tau}, \ref{s2:lemCap}, \ref{s2:lemHS}, \ref{s1:citDef11}, \ref{s1:citLem14}, \ref{s2:lemSEP}, \ref{s1:condG1}}
\end{proposition}

\begin{proof}
By \Gc{1}(a) (Condition~\ref{s1:condG1}) at $\mu=\log\log\Dstar$, we have
$\log\log\Dstar\ge2^8$, so $\lambda_l\ge\log\Dstar\ge2^{256}$ for every $l\le R$;
we use only much weaker consequences, such as $P_l\ge11$ and $\Dstar\ge8$.

(iv) Let $v$ lie in the light part of a light pre-part $Z$ of round $l$, i.e.\
$v\in Z^0\setminus S_Z$. If $v\in D_l$ then, as $v\notin S_Z$,
$\home_l(v)=Z$. If $v\notin D_l$ then $Z$ is the only round-$l$ pre-part
containing $v$, so again $\home_l(v)=Z$. Hence $v$ lies only in the light part
of $\home_l(v)$ (if that pre-part is light), which proves the first and last
assertions. $|Y|=|Y^0|-|S_Y|\ge|Y^0|/2$ by (L1), and $|Y^0|\ge P_r$. By
definition of $D_l$, a vertex outside $D_l$ lies in at most one round-$l$
pre-part, so the sets $U_Z=Z^0\setminus D_l$ are pairwise disjoint.

(i) Each $X^0_Z$ is a leaf of a $\tau$-run, hence a $(2^{-5},s_l)$-expander by
the stopping rule (Lemma~\ref{s2:lem14tau}(a)). For a light pre-part $Z$ apply
Lemma~\ref{s2:lemHS} to $X=X^0_Z$, $z=|Z^0|\ge P_l\ge11$, $W=S_Z$ and
$s_W=s_l/2$: $|S_Z|\le|Z^0|/2$ by (L1), and $X^0_Z$ is a subgraph of
$G'_l[Z^0]$ (node graphs of a split recursion are subgraphs of the root), so by
(L2) every vertex has at most $s_l/2$ neighbours in $S_Z$ in $X^0_Z$. So
$X_Z=X^0_Z-S_Z$ is a $(2^{-6},s_l/2)$-expander on $Z^0\setminus S_Z$. For an
ancestor $Y$, $H_Y$ is an $(\eps_Y,s_Y)$-expander on $|V(Y)|\ge P_r/2\ge2$
vertices (by (iv)), so its minimum degree exceeds $s_Y$ by the remark in Cited
result~\ref{s1:citDef11}. Finally $M_r\ge2^{16}T\log^4T\ge T=d_r\log^2d_r\ge d_r$
(with $T=\tHB_rd_r$), so $\Lambda_r\ge\lambda_r$ and
$s_r\ge\Lambda_r^{100}\ge\lambda_r^{100}$.

(ii) is Lemma~\ref{s2:lemCap}(ii).

(iii) By (R1), $E(G_l)=E(\Cyc_l)\sqcup E(G'_l)$. By Lemma~\ref{s2:lemSEP}(i)
applied to the two-level recursion of round $l$, every edge of $G'_l$ is
deleted exactly once or lies in exactly one leaf. By (R5), every edge of $G'_l$
is assigned to exactly one light part or standalone pre-part, or passes to
$G_{l+1}$: step (1) assigns each edge at most once (the $X^0_Z$ are
edge-disjoint and $X_Z\subseteq X^0_Z$), and steps (2)--(4) treat each
remaining edge once. So $E(G'_l)=\bigsqcup_{Z\text{ light}}E_l(Z)\sqcup
\bigsqcup_{Z\in\Std_l}E_l(Z)\sqcup E(G_{l+1})$; in particular
$E(G_{l+1})\subseteq E(G'_l)\subseteq E(G_l)$. By (R0), $E(G_{R+1})=E_0$, and
iterating from $l=1$ to $l=R$ gives the displayed partition. By (R5)(1),
$E(X_Y)$ is assigned to a light $Y$ and $E(X^0_Y)$ to a standalone $Y$, so
$E(H_Y)\subseteq E_{r(Y)}(Y)$; and each step of (R5) assigns an edge to $Y$
only if both its ends lie in $V(Y)$. Two pre-parts of the same round are
distinct leaves of the two-level recursion, so their leaf graphs are
edge-disjoint by Lemma~\ref{s2:lemSEP}(i). For distinct ancestors $Y\ne Y'$ (of
any rounds), $E(H_Y)\subseteq E_{r(Y)}(Y)$ and $E(H_{Y'})\subseteq
E_{r(Y')}(Y')$, and these assigned sets are disjoint by the partition above.
(The graphs $X^0_Y$ of light pre-parts of different rounds need not be
edge-disjoint: an edge of $X^0_Y$ between a guest $x\in S_Y$ and the light part
$Y$ may be assigned by no step of (R5) at round $r(Y)$, pass down, and lie in a
leaf graph of a later round. Only the statement above is claimed.)

\emph{Long cycles.} For $l\le R$ put $y_l\defeq d_l\ge\Dstar$; by the first part
of (iii), $y_{l+1}\le y_l$. The cycles of $\Cyc_l$ are edge-disjoint, each has
length at least $\tHB_ld_l=y_l\log^2y_l$, and together they have at most
$|E(G_l)|-|E(G'_l)|\le|E(G_l)|-|E(G_{l+1})|=\frac n2(y_l-y_{l+1})$ edges. Hence
\[
|\Cyc_l|\le\frac n2\cdot\frac{y_l-y_{l+1}}{y_l\log^2y_l}.
\]
Group the rounds $l\le R$ by $k\defeq\lfloor\log y_l\rfloor\ge k_0\defeq
\lfloor\log\Dstar\rfloor$. The rounds with a given $k$ are consecutive, and for
them $y_l\ge2^k$ and $\log^2y_l\ge k^2$, while the sum of $y_l-y_{l+1}$ over
them is at most the first $y_l$ of the group, which is less than $2^{k+1}$. So
these rounds contribute at most $\frac n2\cdot\frac{2^{k+1}}{2^kk^2}=n/k^2$,
and $\sum_{l\le R}|\Cyc_l|\le n\sum_{k\ge k_0}k^{-2}\le n/(k_0-1)\le n$, as
$k_0\ge3$.

Finally $E_0=E(G_{R+1})$ and $d_{R+1}<\Dstar$ by (R0), so
$|E_0|=nd_{R+1}/2<\Dstar n/2$.
\end{proof}

\begin{lemma}[edge laminarity, EL]\label{s2:lemEL}
For every ancestor $Y$ of round $r$ and every round $l>r$, no edge of $G_l$ has
both ends in $V(Y)$.
\deps{\ref{s2:defHBtp}, \ref{s2:defAncestors}}
\end{lemma}

\begin{proof}
By (R1) and (R5)(4), $E(G_{j+1})\subseteq E(G'_j)\subseteq E(G_j)$ for every
round $j$, so $E(G_l)\subseteq E(G_{r+1})\subseteq E(G'_r)$. Let $e\in
E(G_{r+1})$; then $e$ is an edge of $G'_r$ that was not assigned at round $r$.
If $Y$ is light and both ends of $e$ lie in $V(Y)=Y$, then $e$, if not assigned
in step (1) of (R5), is assigned in step (2) to the first light part containing
both ends; a contradiction. If $Y$ is standalone (a GC-part or not) and both
ends lie in $V(Y)=Y^0$, then $e$, if not assigned in steps (1) or (2), is
assigned in step (3) to the first standalone pre-part whose vertex set contains
both ends; again a contradiction.
\end{proof}

\begin{proposition}[overlap constants on $\HBtp$]\label{s2:propOV}
Fix a valid $\HBtp$ run and a round $r\le R$. Let $S_r$ be the total size of the
leaves of the two-level recursion of round $r$, and $S^{\piece}_r$ the total
size of the $s=0$ pieces. Then $S^{\piece}_r\le1.12\,n$ and $S_r\le1.21\,n$,
and:
\begin{enumerate}
\item[(K1)]\namedlabel{s2:K1}{(K1)} $\sum|Z^0|\le1.37\,n$, the sum over the
round-$r$ pre-parts; $|\Std_r|\le1.37\,n/P_r$; the number of ancestors of round
$r$ is at most $1.37\,n/P_r$ (each ancestor corresponds to a distinct pre-part
$Y^0\supseteq V(Y)$); hence the number $\nu_l$ of ancestors of rounds at most
$l-2$ satisfies $\nu_l\le1.37\,n\sum_{r\le l-2}P_r^{-1}$;
\item[(K2)]\namedlabel{s2:K2}{(K2)} $|D_r|\le\dup_r\le7.6\,\eps n/\log P_r$,
and $\sum_{Z\in\Std_r}|A_Z|\le\sum|Z^0\cap D_r|\le2\dup_r\le15.2\,\eps n/\log
P_r$, the middle sum over all round-$r$ pre-parts;
\item[(K3)]\namedlabel{s2:K3}{(K3)} the total size $\sum|Z^0|$ of the round-$r$
pre-parts failing (L1) is at most $30.4\,\eps n/\log P_r$, and
$\sum_Y|Y^0\cap\Dups_r|\le16\,\eps n/\log P_r$, the sum over all round-$r$
pre-parts $Y$;
\item[(F11)] $\sum_w\mult_r(w)=\sum_Y|V(Y)|\le1.37\,n$, the second sum over the
ancestors $Y$ of round $r$, and $\mult_r(w)\le\mu_r(w)$ for every vertex $w$.
\end{enumerate}
\deps{\ref{s2:lemOVgeneric}, \ref{s2:lem14tau}, \ref{s2:propStructure}, \ref{s2:defAncestors}, \ref{s2:defHBtp}, \ref{s2:lemSEP}}
\fixes{\RTb{I8}}
\end{proposition}

\begin{proof}
\emph{The composite tree.} The two-level recursion of round $r$ is a split
recursion on $G'_r$, whose vertex set $V(G)$ has $n$ elements. Its splits in the
first level are splits of an $s=0$ recursion, which satisfy the hypotheses of
Lemma~\ref{s2:lemOVgeneric} with $c=1$ (the instance $c=1$ there), hence with
$c=1.6$. Its splits in the second level are splits of the $\tau$-run of a piece
$\piece$ with $|\piece|\ge P_r$; since $\tau_r\ge128\,s_r\log^2|\piece|$
(Proposition~\ref{s2:propStructure}(ii)), Lemma~\ref{s2:lem14tau}(b) shows that
they satisfy the hypotheses with $c=1.6$. Every non-leaf node has at least two
vertices. So Lemma~\ref{s2:lemOVgeneric} applies to the whole two-level
recursion with $c=1.6$ and $n_0=n$: $S_r\le n/(1-5.472\eps)\le1.21\,n$, and for
$M\ge2$,
\begin{equation}\label{s2:eqDupComposite}
\Delta_{\ge M}\le\frac{5.46\,\eps S_r}{\log M}\le\frac{6.61\,\eps n}{\log M}.
\end{equation}
Applied to the first level alone (an $s=0$ recursion on $n$ vertices, $c=1$),
the lemma gives $S^{\piece}_r\le1.12\,n$.

(K1) Pre-parts are leaves, so $\sum|Z^0|\le S_r\le1.21\,n\le1.37\,n$. Every
pre-part has at least $P_r$ vertices, so there are at most $1.37\,n/P_r$
pre-parts. $\Std_r$ is a set of pre-parts, and the map from ancestors of round
$r$ to pre-parts (a light part to its pre-part, a standalone pre-part to
itself) is injective. Summing over $r\le l-2$ gives the bound on $\nu_l$.

(K2) The round-$r$ pre-parts are exactly the leaves of the two-level recursion
with at least $P_r$ vertices (by (R3) its leaves are the pieces with fewer than
$P_r$ vertices and the leaves of the $\tau$-runs; a piece with at least $P_r$
vertices whose $\tau$-run is a single node is such a leaf, hence a pre-part). By
Lemma~\ref{s2:lemOVgeneric}(b) with $M=P_r\ge2$, $\mu_r(w)\le1+\dup_{\ge
P_r}(w)$ for every $w$. Hence $\dup_r=\sum_w(\mu_r(w)-1)^+\le\Delta_{\ge P_r}\le
6.61\,\eps n/\log P_r\le7.6\,\eps n/\log P_r$ by \eqref{s2:eqDupComposite}.
Every $w\in D_r$ contributes at least $1$ to $\dup_r$, so $|D_r|\le\dup_r$.
Next, $\sum_{Z\in\Std_r}|A_Z|=\sum_{Z\in\Std_r}|Z^0\cap D_r|\le\sum_Z|Z^0\cap
D_r|=\sum_{w\in D_r}\mu_r(w)\le\sum_{w\in D_r}2(\mu_r(w)-1)=2\dup_r$, using
$\mu\le2(\mu-1)$ for $\mu\ge2$.

(K3) If $Z$ fails (L1) then $|Z^0|<2|S_Z|\le2|Z^0\cap D_r|$; summing and using
(K2) gives at most $4\dup_r\le30.4\,\eps n/\log P_r$. For the second bound, let
$u\in\Dups_r$ lie in a pre-part $Y$. Then $u$ lies in another leaf $\Leaf\ne Y$
of the two-level recursion. Let $\nu$ be the lowest common ancestor of the
leaves $Y$ and $\Leaf$; they lie below different children of $\nu$, and $u$
lies in both (vertex sets shrink downwards), so $u\in N''_\nu$ by
Lemma~\ref{s2:lemSEP}(0); and $|\nu|\ge|Y^0|\ge P_r$. Hence
$\dup_{\ge P_r}(u)\ge1$ and $\mu_r(u)\le1+\dup_{\ge P_r}(u)\le2\dup_{\ge
P_r}(u)$. Summing over $u\in\Dups_r$,
\[
\sum_Y|Y^0\cap\Dups_r|=\sum_{u\in\Dups_r}\mu_r(u)\le2\Delta_{\ge P_r}
\le\frac{13.3\,\eps n}{\log P_r}\le\frac{16\,\eps n}{\log P_r}.
\]

(F11) Counting pairs $(w,Y)$ with $w\in V(Y)$ gives
$\sum_w\mult_r(w)=\sum_Y|V(Y)|$. The injective map $Y\mapsto Y^0$ of (K1)
satisfies $V(Y)\subseteq Y^0$, so $\sum_Y|V(Y)|\le\sum_{Z}|Z^0|\le1.37\,n$ and,
restricted to the ancestors containing $w$, it shows $\mult_r(w)\le\mu_r(w)$.
\end{proof}

\subsection{Degree recursion, existence, and tower facts}

\begin{proposition}[degree recursion]\label{s2:propDegRec}
Let $l$ be a round of a valid $\HBtp$ run with $d_l\ge\Dstar$ (that is,
$l\le R$), and put $x\defeq\lambda_l$. Then $M_l\le d_l^2$, $\Lambda_l\le2x$,
$s_l\le P_l$, and
\[
d_{l+1}\le1.21\,P_l+9\,s_l\log M_l\le6P_l+9\,s_l\log M_l+1.2\,s_l\le
7\,x^{103}\le x^{\Aexp}<d_l
\]
($\Aexp=104$ would suffice; $\Aexp=105$ is used). More precisely, every edge of
$G_{l+1}$ is of one of the following kinds:
\begin{enumerate}
\item[(a)] an edge deleted by the $\tau$-run of a piece $\piece$ with
$|\piece|\ge P_l$; there are at most $4s_l|\piece|\log|\piece|$ of them for
each piece and at most $4.5\,n\,s_l\log M_l$ in total;
\item[(b)] an edge of a leaf of the two-level recursion with fewer than $P_l$
vertices (a small leaf of a $\tau$-run, or a piece with $|\piece|<P_l$); a leaf
with $q$ vertices has fewer than $qP_l/2$ edges;
\item[(c)] a \emph{guest edge}: an edge of $X^0_Z$ with an end in $S_Z$, for a
light pre-part $Z$; there are at most $|Z^0|s_l/2$ of them for each light $Z$.
\end{enumerate}
Standalone pre-parts, GC-parts included, pass none of the edges of their leaf
graphs down; in particular, declaring a pre-part a GC-part (instead of light)
can only decrease the set of passed-down edges.
\deps{\ref{s2:lem14tau}, \ref{s2:lemCap}, \ref{s2:propOV}, \ref{s2:defHBtp}, \ref{s2:lemSEP}, \ref{s1:condG1}, \ref{s2:lemOVgeneric}}
\end{proposition}

\begin{proof}
Put $\mu\defeq\log x$, where $x=\log d_l\ge\log\Dstar$; so $\mu\ge\log\log\Dstar$,
and \Gc{1}(a),(b) (Condition~\ref{s1:condG1}) hold at $\mu$: $\mu\ge2^8$ and
$x=2^\mu\ge2^{14}\Aexp\mu^3$. Hence $x\ge2^{256}$, so $x\ge7$, $x^2\ge2^{107}$ and
$x^3\ge2^{101}$; and $x\ge2^{14}\cdot105\log^3x$, so $x\ge21+6\log x$ and
$x>105\log x$.

\emph{Parameters.} Let $T\defeq\tHB_ld_l=2^xx^2$. Then $\log T=x+2\log x\le2x$,
so $2^{16}T\log^4T\le2^{16}\cdot2^xx^2\cdot16x^4=2^{20}x^62^x\le2^{2x-1}$ because
$x\ge21+6\log x$; also $2^{40}\le2^{2x-1}$. By (R2), $M_l$ is at most this
maximum plus $1$, so $M_l\le2^{2x-1}+1\le2^{2x}=d_l^2$ and
$\Lambda_l=\log M_l\le2x$. Then $s_l\le\Lambda_l^{100}+1\le2^{100}x^{100}+1\le
2^{101}x^{100}\le x^{103}\le P_l$.

\emph{Classification.} Let $e\in E(G_{l+1})$; by (R5), $e$ is an edge of $G'_l$
not assigned at round $l$. By Lemma~\ref{s2:lemSEP}(i) applied to the two-level
recursion, $e$ is deleted at some split, or lies in exactly one leaf $\Leaf$.
Splits of the first level delete nothing (Lemma~\ref{s2:lemOVgeneric}, instance
$c=1$), so a deleted $e$ is of kind (a). If $|\Leaf|<P_l$, $e$ is of kind (b).
Otherwise $\Leaf$ is a pre-part $Z$ and $e\in E(X^0_Z)$. If $Z$ is standalone,
$e$ is assigned to $Z$ by (R5)(1), which is impossible. If $Z$ is light and $e$
has no end in $S_Z$, then $e\in E(X_Z)$ is assigned by (R5)(1), again
impossible. So $e$ is of kind (c). The same argument shows the assertion about
standalone pre-parts and GC-parts.

\emph{Counting.} (a) By Lemma~\ref{s2:lemCap}(ii), Lemma~\ref{s2:lem14tau}(a)
applies to the $\tau$-run of $\piece$ and gives at most
$4s_l|\piece|\log|\piece|\le4s_l|\piece|\log M_l$ deleted edges. By
Proposition~\ref{s2:propOV}, the pieces have total size
$S^{\piece}_l\le1.12\,n$, so kind (a) has at most $4.48\,n\,s_l\log M_l$ edges.
(b) A leaf with $q<P_l$ vertices has at most $q(q-1)/2<qP_l/2$ edges.
(c) Every guest edge of $X^0_Z$ joins some $u\in Z^0$ to a vertex of $S_Z$, so
it is counted in $\sum_{u\in Z^0}|N_{X^0_Z}(u)\cap S_Z|$; each term is at most
$s_l/2$ by (L2), since $X^0_Z\subseteq G'_l[Z^0]$. The leaves of kind (b) and
the light pre-parts are distinct leaves of the two-level recursion, so the sum
of $q$ over the leaves of kind (b) plus the sum of $|Z^0|$ over the light
pre-parts is at most $S_l\le1.21\,n$ (Proposition~\ref{s2:propOV}). As
$s_l\le P_l$, kinds (b) and (c) together have at most
$\frac12\cdot1.21\,n\cdot P_l$ edges. Therefore
$|E(G_{l+1})|\le4.5\,n\,s_l\log M_l+0.605\,n\,P_l$, that is,
$d_{l+1}\le9\,s_l\log M_l+1.21\,P_l$, which is at most
$6P_l+9s_l\log M_l+1.2s_l$.

\emph{Final bounds.} Using $P_l\le x^{103}+1$, $s_l\le2^{101}x^{100}$ and
$\log M_l\le2x$:
\[
6P_l+9s_l\log M_l+1.2s_l\le6x^{103}+6+9\cdot2^{102}x^{101}+1.2\cdot2^{101}x^{100}
\le6x^{103}+2^{107}x^{101}\le7x^{103},
\]
since $x^2\ge2^{107}$. Next $7x^{103}\le x^{104}\le x^{105}$ as $x\ge7$, and
$x^{105}<2^x=d_l$ as $105\log x<x$.
\end{proof}

\begin{proposition}[existence and termination]\label{s2:propExists}
For every graph $G$ a valid $\HBtp$ run exists, and every sequence of admissible
choices terminates. More precisely: within a round, (R1) and the first level of
(R3) terminate; the $\tau$-runs terminate with any witnesses ($U'\ne\emptyset$
and $n_1,n_2<m$ at every split); (GC) is a deterministic relabelling that feeds
back into neither $D_l$ nor $\home$ nor the guest sets; and $d_{l+1}<d_l$ while
$d_l\ge\Dstar$, so $R$ is finite.
\deps{\ref{s2:lem14tau}, \ref{s2:propDegRec}, \ref{s2:defHBtp}, \ref{s2:lemCap}, \ref{s2:lemOVgeneric}, \ref{s2:defWitness}, \ref{s2:propOV}}
\end{proposition}

\begin{proof}
The proofs of Lemma~\ref{s2:lemCap}(ii), Proposition~\ref{s2:propOV} and
Proposition~\ref{s2:propDegRec} use only the data of the round $l$ in question
and $d_l\ge\Dstar$; so they apply to every round $l$ with $d_l\ge\Dstar$ of any
(possibly unfinished) execution of the procedure.

Fix a round $l$ with $d_l\ge\Dstar$ (if $d_1<\Dstar$ the run stops at once
with $R=0$). (R1) deletes at least one edge per step, so it terminates. In the
first level of (R3), a node that is not an $(\eps,0)$-expander has a witness
(Definition~\ref{s2:defWitness}), and its children have $|U|+|N|<\frac34m$ and
$m-|U|<m$ vertices (Lemma~\ref{s2:lemOVgeneric}, instance $c=1$); so every
root-to-node path has at most $n$ nodes and this level terminates. By
Lemma~\ref{s2:lemCap}(ii), $\tau_l\ge128\,s_l\log^2|\piece|$ for every piece,
so Lemma~\ref{s2:lem14tau}(a) applies to every $\tau$-run with every choice of
witnesses: $U'\ne\emptyset$, $n_1,n_2<m$, and the $\tau$-run terminates.
In (R4), $D_l$, $\home$ and the sets $S_Z$ are determined by the pre-parts and
the home order alone; conditions (L1), (L2) and (GC) are evaluated afterwards,
and their outcome is not used in defining $D_l$, $\home$ or $S_Z$. (R5) is
deterministic. All required choices (witnesses, vertex orders, home orders)
exist. So each round is a finite procedure. By
Proposition~\ref{s2:propDegRec}, $d_{l+1}<d_l$ whenever $d_l\ge\Dstar$; as
$|E(G_l)|$ is a non-negative integer, (R0) stops after finitely many rounds.
\end{proof}

\begin{lemma}[tower-lacunary sums]\label{s2:lemLacunary}
Let $x_0\defeq\log\Dstar$.
\begin{enumerate}
\item[(i)] If $X_1,\dots,X_R\ge0$ and $X_r\le X_{r+1}/2$ for all $r<R$, then
$\sum_{r\le R}X_r\le2X_R$ and $\sum_{r\le R}(R-r+1)X_r\le4X_R$.
\item[(ii)] For a valid run with $R\ge1$, the degree recursion gives
$\lambda_r\ge2^{\lambda_{r+1}/\Aexp}$ for all $r<R$. Hence, if
$F:[x_0,\infty)\to[0,\infty)$ is non-increasing and $F(2^{x/\Aexp})\le F(x)/2$
for all $x\ge x_0$, then $\sum_{r\le R}F(\lambda_r)\le2F(\lambda_R)\le
2F(x_0)$ and $\sum_{r\le R}(R-r+1)F(\lambda_r)\le4F(x_0)$; for the shifted
sums (when $R\ge3$), $\sum_{l=3}^RF(\lambda_{l-2})\le2F(\lambda_{R-2})\le2F(x_0)$
and $\sum_{l=3}^R(R-l+1)F(\lambda_{l-2})\le4F(\lambda_{R-2})\le4F(x_0)$.
\item[(iii)] More generally, if $F:[x_0,\infty)\to[0,\infty)$ satisfies
\[
\text{(H)}\qquad F(y)\le F(x)/2\quad\text{whenever } x\ge x_0\text{ and }
y\ge2^{x/\Aexp},
\]
then $\sum_{r\le R}F(\lambda_r)\le2F(\lambda_R)$ and $\sum_{r\le
R}(R-r+1)F(\lambda_r)\le4F(\lambda_R)$ (and likewise for $F(\lambda_{l-2})$).
\item[(iv)] Under \Gc{1}, condition (H) holds for $F(x)=x^{-a}$ (every $a\ge
1/100$), for $F(x)=1/\log x$ and for $F(x)=(95\log x+8)\,x^{-b}$ (every $b\ge
1$; not used in this manuscript), all three non-increasing on $[x_0,\infty)$;
and for the three
functions $F(x)=(2\logs(2^x)+2)/\eta(x)$ with $\eta(x)\in\{x,\ x^{1/2},\
\log x\}$ (note $2\logs(2^x)+2=2\logs x+4$ for $x>1$; not used in this
manuscript). Each of these three
satisfies $F(x)\le2F(x_0)$ for all $x\ge x_0$ (not used in this manuscript).
\end{enumerate}
Round-$l$ quantities such as $1/M_l$ or
$M_l^3(95\log\lambda_{l-2}+8)/\lambda_{l-2}^{95}$ are summed either by (i),
once the halving of consecutive terms is known, or by bounding them by
$F(\lambda_{l-2})$ for an $F$ as in (ii); for instance the bound
$M_l\le\lambda_{l-2}^{1.6}$, proved in part~(b) of the next lemma (this
example is not used in the proof of the present lemma), gives
$M_l^3(95\log\lambda_{l-2}+8)/\lambda_{l-2}^{95}\le(95\log\lambda_{l-2}+8)\,
\lambda_{l-2}^{-90.2}$. (This example is not used in this manuscript:
Lemma~\ref{s7:lemUHsplit}(iv) bounds its round-$l$ terms by a function of its
own and sums them by~(i).)
\deps{\ref{s2:propDegRec}, \ref{s1:condG1}}
\end{lemma}

\begin{proof}
(i) By induction $X_r\le2^{-(R-r)}X_R$, so $\sum_rX_r\le X_R\sum_{k\ge0}2^{-k}
=2X_R$ and $\sum_r(R-r+1)X_r\le X_R\sum_{k\ge0}(k+1)2^{-k}=4X_R$.

(iii) For $r<R$, Proposition~\ref{s2:propDegRec} at round $r$ gives
$d_{r+1}\le\lambda_r^{\Aexp}$, i.e.\ $\lambda_{r+1}\le\Aexp\log\lambda_r$,
i.e.\ $\lambda_r\ge2^{\lambda_{r+1}/\Aexp}$; and $\lambda_{r+1}\ge x_0$. By (H)
with $x=\lambda_{r+1}$ and $y=\lambda_r$, $F(\lambda_r)\le F(\lambda_{r+1})/2$,
and (i) applies to $X_r\defeq F(\lambda_r)$. For sums of $F(\lambda_{l-2})$
over $3\le l\le R$ apply the same to $r=l-2\le R-2$.

(ii) Let $x\ge x_0$. Then $\mu\defeq\log x\ge\log x_0=\log\log\Dstar$, so \Gc{1}(a),(b)
at $\mu$ give $\mu\ge2^8$ and $x=2^\mu\ge2^{14}\Aexp\mu^3\ge\Aexp\mu=\Aexp\log x$, that is,
$2^{x/\Aexp}\ge x$. Hence $2^{x/\Aexp}\ge x\ge x_0$ for $x\ge x_0$.
If $F$ is non-increasing and $F(2^{x/\Aexp})\le F(x)/2$, then for $y\ge
2^{x/\Aexp}$ we get $F(y)\le F(2^{x/\Aexp})\le F(x)/2$, which is (H). So (iii)
applies, and $F(\lambda_R)\le F(x_0)$ as $\lambda_R\ge x_0$.

(iv) Let $x\ge x_0$ and $y\ge2^{x/\Aexp}$, and put $\mu\defeq\log x\ge\log\log\Dstar$.
By \Gc{1}(a),(b) at $\mu$, $\mu\ge2^8$ and $x=2^\mu\ge2^{14}\Aexp\mu^3$; hence
$x\ge2^{256}$, $x/\Aexp\ge101\log x$, $x\ge2\Aexp\log x$, $x\ge\Aexp\log^2x/4$ and
$x/\Aexp\ge2^{12}$, and all of this holds in particular for $x=x_0$. Then
$\log y\ge x/\Aexp$, and $2^{x/\Aexp}\ge x^{101}$ because $x/\Aexp\ge101\log x$; so
$y\ge x^{101}\ge2^{32}$.
We use two facts about $\logs$: $\logs z\le z$ for $z\ge1$ (if $z\le2$ then
$\logs z\le1$; if $z>2$ then $\logs z=1+\logs(\log z)\le1+\log z\le z$ by
induction), and hence, for $y\ge4$, $\logs y=2+\logs(\log\log y)\le2+\log\log y
\le1+\log y$.

\emph{$F(x)=x^{-a}$, $a\ge1/100$:} $F(y)/F(x)=(x/y)^a\le x^{-100a}\le x^{-1}\le
1/2$.

\emph{$F(x)=1/\log x$:} $F(y)=1/\log y\le\Aexp/x\le1/(2\log x)=F(x)/2$, as
$x\ge2\Aexp\log x$.

\emph{$F(x)=(95\log x+8)x^{-b}$, $b\ge1$:} $x^{b+1}F'(x)=95/\ln2-b(95\log x+8)
<0$, so $F$ is decreasing; and $F(y)\le y^{-b/2}$, because $95\log y+8\le
y^{1/2}$ for $y\ge2^{30}$ (true at $2^{30}$, and the right side grows faster).
Hence $F(y)\le x^{-50b}\le x^{-b}/2\le F(x)/2$.

\emph{$F(x)=(2\logs x+4)/\eta(x)$.} Here $F(x)\ge4/\eta(x)$, so it suffices to
show $F(y)\le2/\eta(x)$. For $\eta(x)=x$: $F(y)\le(2\log y+6)/y\le y^{-1/2}$
(the inequality $2\log y+6\le y^{1/2}$ holds at $y=2^{12}$, and the right side
grows faster beyond), and $y^{1/2}\ge x^{50}\ge x/2$. For $\eta(x)=x^{1/2}$:
$F(y)\le(2\log y+6)/y^{1/2}\le y^{-1/4}$ (valid for $y\ge2^{32}$ in the same
way), and $y^{1/4}\ge x^{25}\ge x^{1/2}/2$. For $\eta(x)=\log x$: with
$u\defeq\log y\ge x/\Aexp$, $F(y)\le(2\log u+8)/u\le u^{-1/2}$ (valid for $u\ge
2^{12}$), and $u^{-1/2}\le(\Aexp/x)^{1/2}\le2/\log x$ as $x\ge\Aexp\log^2x/4$.

\emph{The bound $F(x)\le2F(x_0)$ for the three $\logs$-functions.} Let
$\mathrm{tw}(0)\defeq1$ and $\mathrm{tw}(j+1)\defeq2^{\mathrm{tw}(j)}$. Since
$\log^{[i]}$ is increasing and $\log^{[j]}\mathrm{tw}(j)=1$, a real $z>1$ has
$\logs z=j$ iff $\mathrm{tw}(j-1)<z\le\mathrm{tw}(j)$. Let $k_0\defeq\logs x_0\ge
1$, let $x\ge x_0$, and write $\logs x=k_0+k$ with $k\ge0$. Then $x_0\le
\mathrm{tw}(k_0)$ and $x>\mathrm{tw}(k_0+k-1)$. If $k=0$, then $F(x)\le F(x_0)$
since $\eta$ is increasing. If $k=1$, then $2\logs x+4=2k_0+6\le2(2k_0+4)$ and
$\eta(x)\ge\eta(x_0)$, so $F(x)\le2F(x_0)$. Let $k\ge2$. We claim
$\eta(x)\ge2^k\eta(x_0)$. Note $\mathrm{tw}(k_0+1)=2^{\mathrm{tw}(k_0)}\ge2^{x_0}$,
and $\mathrm{tw}(j+1)=2^{\mathrm{tw}(j)}\ge4\,\mathrm{tw}(j)$ whenever
$\mathrm{tw}(j)\ge4$. Hence $x>\mathrm{tw}(k_0+k-1)\ge4^{k-2}2^{x_0}$, which gives
the claim for $\eta(x)=x$ (as $4^{k-2}2^{x_0}\ge2^kx_0$ for $k\ge2$) and for
$\eta(x)=x^{1/2}$ (as $2^{k-2}2^{x_0/2}\ge2^kx_0^{1/2}$). For
$\eta=\log$: $\log x>\mathrm{tw}(k_0+k-2)$, which is at least $\mathrm{tw}(k_0)\ge
x_0\ge4\log x_0$ if $k=2$, and at least $4^{k-3}2^{x_0}\ge2^k\log x_0$ if
$k\ge3$. Using the claim and $2k_0+2k+4\le2^k(2k_0+4)$ (as $2^k\ge1+k$),
$F(x)\le(2k_0+2k+4)/(2^k\eta(x_0))\le(2k_0+4)/\eta(x_0)=F(x_0)$.
\end{proof}

\begin{lemma}[tower facts]\label{s2:lemTower}
For every valid $\HBtp$ run with $d_1\ge\Dstar$ (so $R\ge1$):
\begin{enumerate}
\item[(a)] $\lambda_r\ge d_{r+1}^{1/\Aexp}$ for $r\le R$; $d_{r+2}\le\lambda_r$
for $r\le R-1$; $R-r\le2\logs d_r-1$, in particular
\[
R-r\le2\logs d_r+2\qquad(r\le R);
\]
and $\lambda_r\ge\lambda_{l-2}$ whenever $r\le l-2\le R$.
\item[(b)] For $r\le R$: $M_r\le d_r^2$; $\lambda_r\le\Lambda_r\le2\lambda_r$;
$\lambda_r^{100}\le s_r\le2\Lambda_r^{\sigma}$; $s_r\le P_r$; and $L_Y\le\log
M_r\le2\lambda_r$ for every ancestor $Y$ of round $r$. For $3\le l\le R$:
$M_l\le(\Aexp\log\lambda_{l-2})^{2\Aexp}\le\lambda_{l-2}^{1.6}$;
$P_{l-2}\ge M_l^{13}$; and $P_{l-2}/2\ge M_l\log^4M_l$. For $2\le l\le R$:
$M_{l-1}\ge2M_l$. For $r<R$: $P_r\ge2P_{r+1}$. Moreover
$\sum_{l\le R}1/M_l\le2/\Dstar$ and $\sum_{l\le R}\psi(M_l)\le2.1\,\psi(\Dstar)$
with $\psi(x)\defeq(6\log x+12)/x$. Finally, completing (K1) of
Proposition~\ref{s2:propOV}, the number $\nu_l$ of ancestors of rounds at most
$l-2$ satisfies $\nu_l\le2.74\,n/P_{l-2}$ for $3\le l\le R$.
\item[(c)] For $r\le R$: $\tau_r\le257\,\Lambda_r^{\sigma+2}\le
2^{\sigma+11}\lambda_r^{\sigma+2}$ and $\tau_r/P_r\le2^{\sigma+11}/\lambda_r$;
and $\thGC_r(Z^0)\ge P_r\lambda_r^{-1/2}>\tau_r$ for every round-$r$ pre-part
$Z$.
\item[(d)] For $r\le l\le R$: $\lambda_r\ge2^{2\logs d_r+2}\ge2^{R-r}\ge
2^{l-r}$.
\item[(e)] Put
\[
\epsA\defeq\frac{31\,\eps}{\Cp\log\log\Dstar},
\]
a function of $\Dstar$ alone (it depends neither on the run nor on $n$). Then
\[
\sum_{l\le R}\frac{15.2\,\eps}{\log P_l}\le\epsA.
\]
Consequently $\sum_{l\le R}|D_l|\le\epsA n$ and $\sum_{l\le
R}\sum_{Z\in\Std_l}|A_Z|\le\epsA n$.
\end{enumerate}
\deps{\ref{s2:propDegRec}, \ref{s2:lemLacunary}, \ref{s2:defHBtp}, \ref{s1:condG1}, \ref{s1:condG2}, \ref{s2:propStructure}, \ref{s2:lemCap}, \ref{s2:propOV}, \ref{s2:defAncestors}}
\fixes{\RTb{I6}, \RTb{I8}}
\end{lemma}

\begin{proof}
Every round $r\le R$ has $d_r\ge\Dstar$, so $\log\lambda_r\ge\log\log\Dstar$, and
the items of \Gc{1} (Condition~\ref{s1:condG1}) hold at $\mu=\log\lambda_r$ for
every $r\le R$; they also hold at $\mu=\log\log d$ for every $d\ge\Dstar$. Proposition~\ref{s2:propDegRec} applies at every round $r\le R$.

(a) Proposition~\ref{s2:propDegRec} at round $r$ gives $d_{r+1}\le
\lambda_r^{\Aexp}$. For $r\le R-1$ it also applies at round $r+1$, so
$d_{r+2}\le\lambda_{r+1}^{\Aexp}$ with $\lambda_{r+1}=\log d_{r+1}\le\Aexp\log
\lambda_r$; hence $d_{r+2}\le(\Aexp\log\lambda_r)^{\Aexp}\le\lambda_r$, because
with $u\defeq\log\lambda_r$ we have $\Aexp\log(\Aexp u)\le0.8u\le u$ by \Gc{1}(c)
at $u$; equivalently, $(\Aexp\log\lambda_r)^{\Aexp}\le\lambda_r$.

Let $k\defeq\logs d_r$; $k\ge1$ as $d_r>1$. Suppose $R\ge r+2k$. We show
$d_{r+2i}\le\log^{[i]}d_r$ for $0\le i\le k$ by induction: for $i<k$ we have
$r+2i\le R-2$, so $d_{r+2i+2}\le\log d_{r+2i}\le\log(\log^{[i]}d_r)$, where
$\log^{[i]}d_r>1$ as $i<\logs d_r$. Hence $d_{r+2k}\le\log^{[k]}d_r\le1<\Dstar$,
contradicting $r+2k\le R$. Therefore $R-r\le2k-1\le2\logs d_r+2$. (The weaker
form with ``$+2$'' is the one used in later sections, \RTb{I6}.) Finally
$d_{j+1}\le d_j$ for all $j\le R$ (Proposition~\ref{s2:propStructure}(iii)), so
$\lambda$ is non-increasing in the round.

(b) $M_r\le d_r^2$, $\Lambda_r\le2\lambda_r$ and $s_r\le P_r$ are in
Proposition~\ref{s2:propDegRec}; $\Lambda_r\ge\lambda_r$ and
$s_r\ge\lambda_r^{100}$ are in (the proof of)
Proposition~\ref{s2:propStructure}(i); and $s_r\le\Lambda_r^\sigma+1\le
2\Lambda_r^\sigma$. For an ancestor $Y$ of round $r$, $|V(Y)|\le|Y^0|\le M_r$
by Lemma~\ref{s2:lemCap}(ii), so $L_Y\le\log M_r=\Lambda_r\le2\lambda_r$.

Let $3\le l\le R$. By Proposition~\ref{s2:propDegRec} at rounds $l-1$ and
$l-2$, $d_l\le\lambda_{l-1}^{\Aexp}$ and $\lambda_{l-1}\le\Aexp\log
\lambda_{l-2}$; with $M_l\le d_l^2$ this gives
$M_l\le(\Aexp\log\lambda_{l-2})^{2\Aexp}$. With $u\defeq\log\lambda_{l-2}$, \Gc{1}(c) at $u$ gives
$2\Aexp\log(\Aexp u)=210\log(105u)\le1.6u$, so
$M_l\le\lambda_{l-2}^{1.6}$. Hence $M_l^{13}\le\lambda_{l-2}^{20.8}\le
\lambda_{l-2}^{103}\le P_{l-2}$. As $\Lambda_l\ge40$, $\log M_l=\Lambda_l\le
2^{\Lambda_l/4}=M_l^{1/4}$, so $M_l\log^4M_l\le M_l^2\le M_l^{13}/2\le
P_{l-2}/2$.

For $2\le l\le R$: $M_{l-1}\ge d_{l-1}$ (as in
Proposition~\ref{s2:propStructure}(i)), and $M_l\le d_l^2\le
\lambda_{l-1}^{2\Aexp}=(\log d_{l-1})^{2\Aexp}$; with $y\defeq\log d_{l-1}$ and $v\defeq\log y\ge\log\log\Dstar$, \Gc{1}(a),(b) at $v$
give $v\ge2^8$ and $y=2^v\ge2^{14}\Aexp v^3\ge1+2\Aexp v$, that is, $2^y\ge2y^{2\Aexp}$; so
$M_{l-1}\ge2M_l$. For $r<R$:
$P_{r+1}\le\lambda_{r+1}^{103}+1\le(\Aexp\log\lambda_r)^{103}+1\le\lambda_r+1\le
\lambda_r^{103}/2\le P_r/2$, because
$(\Aexp\log\lambda_r)^{103}\le(\Aexp\log\lambda_r)^{\Aexp}\le\lambda_r$ (the
first inequality as $\Aexp\log\lambda_r\ge\Aexp\ge1$, since $\lambda_r\ge2$; the
second from the proof of (a)) and $\lambda_r\ge2$.

By $M_{l-1}\ge2M_l$ and Lemma~\ref{s2:lemLacunary}(i) applied to $X_l\defeq
1/M_l$, $\sum_{l\le R}1/M_l\le2/M_R\le2/d_R\le2/\Dstar$. The function $\psi$ is
decreasing on $[1,\infty)$, because $\psi'(x)=(6/\ln2-6\log x-12)/x^2$ and
$6/\ln2<12$; and for $2\le l\le R$,
\[
\psi(M_{l-1})\le\psi(2M_l)=\frac{6\log M_l+18}{2M_l}
=\psi(M_l)\Bigl(\frac12+\frac3{6\log M_l+12}\Bigr)\le0.512\,\psi(M_l),
\]
as $\log M_l\ge40$. Hence $\sum_{l\le R}\psi(M_l)\le\psi(M_R)/(1-0.512)\le
2.05\,\psi(M_R)\le2.1\,\psi(\Dstar)$, since $M_R\ge d_R\ge\Dstar$.

Finally, by (K1) of Proposition~\ref{s2:propOV},
$\nu_l\le1.37\,n\sum_{r\le l-2}1/P_r$, and $P_r\ge2P_{r+1}$ for $r<l-2\le R$;
Lemma~\ref{s2:lemLacunary}(i) applied to $X_r\defeq1/P_r$ ($r\le l-2$) gives
$\sum_{r\le l-2}1/P_r\le2/P_{l-2}$, so $\nu_l\le2.74\,n/P_{l-2}$.

(c) $\tau_r=\lceil128\,s_r\Lambda_r^2\rceil\le128\cdot2\Lambda_r^\sigma
\Lambda_r^2+1\le257\,\Lambda_r^{\sigma+2}\le257\cdot2^{\sigma+2}
\lambda_r^{\sigma+2}<2^{\sigma+11}\lambda_r^{\sigma+2}$, as $257\cdot4<2^{11}$.
Since $\Cp=\sigma+3$, $P_r\ge\lambda_r^{\sigma+3}$ and
$\tau_r/P_r\le2^{\sigma+11}/\lambda_r$. For a round-$r$ pre-part $Z$,
$\thGC_r(Z^0)\ge|Z^0|\lambda_r^{-1/2}\ge P_r\lambda_r^{-1/2}\ge
\lambda_r^{102.5}$, while $\tau_r<2^{111}\lambda_r^{102}$; and
$\lambda_r^{1/2}>2^{111}$ because $\log\lambda_r\ge2^8>222$ by \Gc{1}(a).

(d) Put $\mu\defeq\log\lambda_r$; so $\mu\ge2^8$ by \Gc{1}(a). As $d_r=2^{\lambda_r}$ and
$\lambda_r=2^\mu$ with $\lambda_r,\mu>1$, $\logs d_r=2+\logs\mu$. By the fact
$\logs z\le1+\log z$ ($z\ge4$) shown in the proof of
Lemma~\ref{s2:lemLacunary}(iv), $2\logs d_r+2=6+2\logs\mu\le8+2\log\mu\le\mu$,
the last step by \Gc{1}(d) at $\mu$.
So $\lambda_r=2^\mu\ge2^{2\logs d_r+2}$, which is at least $2^{R-r}$ by (a) and
at least $2^{l-r}$ for $l\le R$.

(e) We have $\log P_l\ge\Cp\log\lambda_l$. The function
$F(x)\defeq1/(\Cp\log x)$ is non-increasing and satisfies
$F(2^{x/\Aexp})\le F(x)/2$ for $x\ge\log\Dstar$
(Lemma~\ref{s2:lemLacunary}(iv)). By Lemma~\ref{s2:lemLacunary}(ii),
\[
\sum_{l\le R}\frac1{\log P_l}\le\sum_{l\le R}F(\lambda_l)\le2F(\log\Dstar)=
\frac2{\Cp\log\log\Dstar},
\]
so $\sum_{l\le R}15.2\,\eps/\log P_l\le30.4\,\eps/(\Cp\log\log\Dstar)\le
\epsA$. By (K2) of Proposition~\ref{s2:propOV},
$\sum_l|D_l|\le\sum_l7.6\,\eps n/\log P_l\le\epsA n$ and
$\sum_l\sum_{Z\in\Std_l}|A_Z|\le\sum_l15.2\,\eps n/\log P_l\le\epsA n$.
\end{proof}

\subsection{Rule (GC), ORIGIN\texorpdfstring{$^\tau$}{-tau}, and parentless mass}

\begin{lemma}[rule GC]\label{s2:lemGC}
For every valid $\HBtp$ run and every round $r\le R$:
\begin{enumerate}
\item[(i)] the pre-parts, their leaf graphs, $\Dups_r$, $D_r$, $\home_r$ and
the guest sets $S_Z$ are defined in (R1)--(R4) before and independently of
condition (GC); (GC) only decides whether a pre-part satisfying (L1) and (L2)
is light or standalone;
\item[(ii)] for every light part $Y$ of round $r$ and every guest $x\in S_Y$,
$e_{X^0_Y}(x,Y)<\thGC_r(Y^0)$, where $Y$ denotes the vertex set
$Y^0\setminus S_Y$;
\item[(iii)] GC-parts are standalone; for a GC-part $Y$, every edge of $G'_r$
with both ends in $Y^0$ is assigned at round $r$ (the edges of $X^0_Y$, those
between a guest and $Y^0\setminus S_Y$ included, by (R5)(1)), so no edge of
$G_{r+1}$ has both ends in $V(Y)=Y^0$; and every guest $x\in S_Y$ lies in
$D_r\subseteq\Dups_r$;
\item[(iv)] $\thGC_r(Z^0)>\tau_r$ for every round-$r$ pre-part $Z$.
\end{enumerate}
\deps{\ref{s2:defHBtp}, \ref{s2:lemTower}, \ref{s2:lemEL}, \ref{s2:defAncestors}}
\end{lemma}

\begin{proof}
(i) is the order of definitions in (R3)--(R4): (GC) is evaluated only after
the pre-parts, $D_r$, $\home_r$ and the sets $S_Z$ are fixed, and none of
these uses the outcome of (GC).

(ii) A light pre-part satisfies (GC): no $x\in S_Y$ has $\thGC_r(Y^0)$ or more
neighbours in $Y^0\setminus S_Y$ in $X^0_Y$.

(iii) A GC-part fails (GC) and is therefore standalone by (R4). By (R5)(1) all
edges of $X^0_Y$ are assigned to $Y$. By Lemma~\ref{s2:lemEL}, applied to the
ancestor $Y$ (with $V(Y)=Y^0$), no edge of $G_{r+1}$ has both ends in $Y^0$;
equivalently, every edge of $G'_r$ inside $Y^0$ is assigned at round $r$. A
guest $x\in S_Y$ lies in $D_r$ by definition, hence in at least two pre-parts,
which are two leaves of the two-level recursion; so $x\in\Dups_r$.

(iv) is Lemma~\ref{s2:lemTower}(c).
\end{proof}

\begin{proposition}[ORIGIN$^\tau$]\label{s2:propOrigin}
Fix a valid $\HBtp$ run, a round $r\le R$, a round-$r$ pre-part $Y$ and a
vertex $u\in Y^0\setminus\Dups_r$.
\begin{enumerate}
\item[(a)] $u$ lies in a unique $s=0$ piece $\piece$ and in a unique leaf of the
two-level recursion, namely $Y$ (so $Y^0\subseteq V(\piece)$). Every edge $ux$
of $G_{r+1}$ is of exactly one of the following two types:
\begin{itemize}
\item[($\beta$)] $x\notin Y^0$, and $ux$ was deleted by the $\tau$-run of
$\piece$ (so $x\in V(\piece)$);
\item[($\alpha$)] $Y$ is light, $x\in S_Y$, and $ux\in E(X^0_Y)$ (a
\emph{guest--core edge}: it is not deleted and not assigned at round $r$, and
passes down).
\end{itemize}
A standalone $Y$ (GC-parts included) has no edges of type ($\alpha$). Since
$E(G_l)\subseteq E(G_{r+1})$ for $l>r$, the same classification applies to the
edges $ux$ of $G_l$.
\item[(b)] For every vertex $h$ and every round $l>r$,
\[
\#\bigl\{u\in Y^0\setminus\Dups_r:\ hu\in E(G_l)\text{ is of type }(\beta)
\bigr\}\le\tau_r-1.
\]
\item[(c)] $\sum_Y|Y^0\cap\Dups_r|\le16\,\eps n/\log P_r$, the sum over all
round-$r$ pre-parts $Y$.
\end{enumerate}
An edge of type ($\alpha$) has its core end $u$ outside $\Dups_r$ (it is not
true that ($\alpha$)-edges occur only at vertices of $\Dups_r$); a guest $x$
of a light $Y$ is the guest end of fewer than $\thGC_r(Y^0)$ edges of type
($\alpha$) (Lemma~\ref{s2:lemGC}(ii)); and there is no third type.
\deps{\ref{s2:lemThinCut}, \ref{s2:lemSEP}, \ref{s2:lemCap}, \ref{s2:propOV}, \ref{s2:lemEL}, \ref{s2:defHBtp}, \ref{s2:lemGC}, \ref{s2:lemOVgeneric}, \ref{s2:propStructure}}
\fixes{\RTb{I5}}
\end{proposition}

\begin{proof}
(a) As $u\notin\Dups_r$, $u$ lies in exactly one leaf of the two-level
recursion, and $u\in Y^0$, so that leaf is $Y$. If $u$ lay in two $s=0$ pieces,
it would lie in a leaf below each of them (Lemma~\ref{s2:lemSEP}(0)), i.e.\ in
two leaves; so $u$ lies in a unique piece $\piece$, and $Y$ is a leaf of the
$\tau$-run of $\piece$ (by Lemma~\ref{s2:lemCap}(ii) this is a
Lemma~14$^\tau$ run, although only its split structure is used below). By
Lemma~\ref{s2:lemSEP}(i),(ii) applied to the two-level recursion, every edge of
$G'_r$ at $u$ is either an edge of the leaf $Y$, i.e.\ of $X^0_Y$, or deleted
at some split on the path from the root to $Y$. The splits of the first level
delete nothing (Lemma~\ref{s2:lemOVgeneric}, instance $c=1$), so deleted edges
at $u$ are deleted by the $\tau$-run of $\piece$ and are edges of $\piece$. A
deleted edge $ux$ has $x\notin Y^0$: at the split where it is deleted, $u$ and
$x$ go to different children, each to only one, so they do not lie in a common
leaf. Hence every edge $ux$ of $G'_r$ with $x\in Y^0$ is an edge of $X^0_Y$.

Now let $ux\in E(G_{r+1})$, an edge of $G'_r$ not assigned at round $r$. Since
$u\notin D_r$ (as $D_r\subseteq\Dups_r$), $Y$ is the only pre-part containing
$u$, and $u\notin S_Y$.

\emph{Case $ux$ deleted.} Then $x\notin Y^0$ and $x\in V(\piece)$; this is type
($\beta$). (Such an edge is indeed never assigned: a light part or standalone
pre-part containing both ends would contain $u$, hence would be the light part
of $Y$ or the pre-part $Y$ itself, neither of which contains $x\notin Y^0$.)

\emph{Case $ux\in E(X^0_Y)$.} If $Y$ is standalone, $ux$ is assigned to $Y$ by
(R5)(1), a contradiction. So $Y$ is light. If $x\notin S_Y$, then $ux\in
E(X_Y)$ is assigned by (R5)(1), a contradiction. So $x\in S_Y$: type
($\alpha$). (Such an edge is indeed passed down: the only light part that could
contain $u$ is $Y$, which does not contain $x$, and the only pre-part
containing $u$ is the light pre-part $Y$, so steps (2) and (3) of (R5) do not
apply.)

The types are exclusive ($x\notin Y^0$ in ($\beta$), $x\in S_Y\subseteq Y^0$ in
($\alpha$)). For standalone $Y$, Lemma~\ref{s2:lemEL} also shows directly that
no edge of $G_{r+1}$ has both ends in $V(Y)=Y^0$. By
Proposition~\ref{s2:propStructure}(iii) (or directly from (R1) and (R5)),
$E(G_l)\subseteq E(G_{r+1})$ for $l>r$.

(b) Let $h$ be a vertex and $l>r$. By (a), each $u$ counted has $hu$ deleted by
the $\tau$-run of the unique piece $\piece$ containing $u$, and this piece is
the piece whose $\tau$-run has $Y$ as a leaf; so all counted edges are deleted
edges of one $\tau$-run, and each is of the form $hu$ with $u\in
Y^0\setminus\Dups_r$. A vertex lying in two leaves of the $\tau$-run of
$\piece$ lies in two leaves of the two-level recursion, so
$Y^0\setminus\Dups_r\subseteq Y^0\setminus\Dup$, where $\Dup$ is taken in the
$\tau$-run of $\piece$. Lemma~\ref{s2:lemThinCut}, applied to this $\tau$-run
(all of whose splits are $\tau$-rule splits at threshold $\tau_r$) and its
leaf $Y$, bounds the number of such edges by $\lceil\tau_r\rceil-1=\tau_r-1$
($\tau_r$ is an integer by (R2)), and by $0$ if $h\in Y^0$. Distinct $u$ give
distinct edges $hu$.

(c) is (K3) of Proposition~\ref{s2:propOV}.

The final paragraph of the statement: an ($\alpha$)-edge is an edge $ux$ with
$u\in Y^0\setminus\Dups_r$ by the setting of (a); the ($\alpha$)-edges at a
guest $x$ with core ends in $Y\setminus\Dups_r$ are edges of $X^0_Y$ between
$x$ and $Y=Y^0\setminus S_Y$, and there are fewer than $\thGC_r(Y^0)$ of them
by Lemma~\ref{s2:lemGC}(ii).
\end{proof}

\begin{proposition}[fresh and parentless mass]\label{s2:propParentless}
For every valid $\HBtp$ run with $d_1\ge\Dstar$:
\begin{enumerate}
\item[(i)] If $x$ lies in a round-$r$ pre-part, then $H\defeq\home_r(x)$ is an
ancestor of round $r$ containing $x$: $x\in H^0\setminus S_H=V(H)$ if $H$ is
light, and $x\in H^0=V(H)$ otherwise. Hence, for $Z\in\Std_l$ and $x\in U_Z$,
$x\in F_Z$ iff $j_0(x)\ge l-1$. A vertex is a port of at most one part per
round, so it is a fresh port only at rounds $j_0(x)$ and $j_0(x)+1$,
and\namedlabel{s2:K4}{(K4)}
\begin{equation*}
\sum_{l\le R}\ \sum_{Z\in\Std_l}|F_Z|\le2n.\tag*{(K4)}
\end{equation*}
\item[(ii)] (admissible parents) If $Y$ is a light part of round $r\le l-2$ and
$Z$ is a light part of round $l\le R$, then
\[
|Y|\ge P_r/2\ge P_{l-2}/2\ge M_l\log^4M_l\ge|Z|\,L_Z^4.
\]
\item[(iii)] Let $\Bad$ be \emph{any} set of light parts. Call a pair $(v,Z)$,
with $Z$ a light part of some round $l$ and $v\in Z$, \emph{parentless with
respect to $\Bad$} if no light part outside $\Bad$ of a round at most $l-2$
contains $v$. Then the number of such pairs is at most
\[
2n+\sum_{Y\in\Bad}|Y|\,\bigl(R-r(Y)\bigr).
\]
More precisely, a vertex $v$ is in parentless pairs only at its light parts of
rounds $j_1(v)$ and $j_1(v)+1$, and at the later rounds at which the light part
of round $j_1(v)$ containing $v$ lies in $\Bad$.
\end{enumerate}
\deps{\ref{s2:propStructure}, \ref{s2:lemTower}, \ref{s2:lemCap}, \ref{s2:defAncestors}, \ref{s2:defHBtp}}
\fixes{\RTb{I9}}
\end{proposition}

\begin{proof}
(i) By definition $H=\home_r(x)$ is a pre-part containing $x$. If $H$ is light,
then $x\notin S_H$, because $S_H$ consists of vertices whose home is not $H$;
so $x\in H^0\setminus S_H=V(H)$, and the light part $H$ is an ancestor of round
$r$. If $H$ is standalone, it is an ancestor with $V(H)=H^0\ni x$.

Let $Z\in\Std_l$ and $x\in U_Z\subseteq Z^0$, so $j_0(x)\le l$. If $l\le2$,
then $x\in F_Z$ by definition and $j_0(x)\ge1\ge l-1$. Let $l\ge3$. If
$j_0(x)\le l-2$, then $\home_{j_0(x)}(x)$ is an ancestor of round $j_0(x)\le
l-2$ containing $x$, so $\anc_l(x)\ne\emptyset$ and $x\notin F_Z$. Conversely,
if $\anc_l(x)\ne\emptyset$, then $x\in V(Y)\subseteq Y^0$ for a pre-part $Y$ of
a round at most $l-2$, so $j_0(x)\le l-2$. Hence $x\in F_Z$ iff $j_0(x)\ge
l-1$.

The port sets $U_Z$ of distinct $Z\in\Std_l$ are disjoint
(Proposition~\ref{s2:propStructure}(iv)), so $x$ is a port of at most one part
per round. If $x$ is a fresh port at round $l$ then $j_0(x)\le l$ (it lies in
$Z^0$) and $j_0(x)\ge l-1$, so $l\in\{j_0(x),j_0(x)+1\}$. Summing over $x$
gives $\sum_l\sum_Z|F_Z|\le2n$.

(ii) $|Y|\ge P_r/2$ by Proposition~\ref{s2:propStructure}(iv). As $\lambda$ is
non-increasing in the round (Lemma~\ref{s2:lemTower}(a)) and $r\le l-2$,
$P_r\ge P_{l-2}$. As $l\ge r+2\ge3$, $P_{l-2}/2\ge M_l\log^4M_l$ by
Lemma~\ref{s2:lemTower}(b). Finally $|Z|\le|Z^0|\le M_l$ by
Lemma~\ref{s2:lemCap}(ii), and $L_Z=\log|Z|\le\log M_l$.

(iii) Fix a vertex $v$ that lies in some light part, and let $j_1\defeq
j_1(v)$ and $Y_1(v)$ be the light part of round $j_1$ containing $v$. By
Proposition~\ref{s2:propStructure}(iv), $v$ lies in at most one light part per
round, so the pairs $(v,Z)$ correspond to distinct rounds $l\ge j_1$. If $l\ge
j_1+2$ and $Y_1(v)\notin\Bad$, then $Y_1(v)$ is a light part outside $\Bad$ of
round $j_1\le l-2$ containing $v$, so $(v,Z)$ is not parentless. Hence the
parentless pairs of $v$ are at the rounds $j_1$ and $j_1+1$ (at most two
pairs), and, if $Y_1(v)\in\Bad$, at rounds $l$ with $j_1+2\le l\le R$ (at most
$R-j_1$ pairs). Summing over $v$, the number of parentless pairs is at most
\[
2n+\sum_{v:\,Y_1(v)\in\Bad}\bigl(R-r(Y_1(v))\bigr)\le2n+\sum_{Y\in\Bad}
\sum_{v\in Y}\bigl(R-r(Y)\bigr)=2n+\sum_{Y\in\Bad}|Y|\bigl(R-r(Y)\bigr),
\]
since $v\in Y_1(v)$.

Part (iii) holds for every set $\Bad$. Its argument uses only that the light
part of round $j_1(v)$ containing $v$ is a light part outside $\Bad$; so, when
it is applied, every reason for which a light part is not used as a parent must
place that part in $\Bad$. In the later application this means that $\Bad$
consists of the light parts that are demoted \emph{or} parent-bad, not only the
demoted ones (\RTb{I9}).
\end{proof}

\begin{table}[ht]
\begin{center}
\begin{tabular}{p{4.6cm}p{2.6cm}p{3.4cm}p{3.2cm}}
\hline
quantity & plain $\HBs$ & $\HBtp$ (used here) & where proved\\
\hline
$\sum|Z^0|$ over round-$r$ pre-parts & $1.28\,n$ & $1.37\,n$ ($1.21\,n$ on
the composite tree) & Prop.~\ref{s2:propOV}, (K1)\\
duplication at nodes of size $\ge P$ & $3.5\,\eps S/\log P$ &
$5.5\,\eps S/\log P$ & Lemma~\ref{s2:lem14tau}(b)\\
ancestors of rounds $\le l-2$ ($\nu_l$) & $2.6\,n/P_{l-2}$ &
$2.74\,n/P_{l-2}$ & Lemma~\ref{s2:lemTower}(b)\\
$\dup_r$ (and $|D_r|$) & $4.5\,\eps n/\log P_r$ & $7.6\,\eps n/\log P_r$ &
Prop.~\ref{s2:propOV}, (K2)\\
$\sum_{Z\in\Std_r}|A_Z|$ & $9\,\eps n/\log P_r$ & $15.2\,\eps n/\log P_r$ &
Prop.~\ref{s2:propOV}, (K2)\\
mass of pre-parts failing (L1) & $18\,\eps n/\log P_r$ &
$30.4\,\eps n/\log P_r$ & Prop.~\ref{s2:propOV}, (K3)\\
\hline
\end{tabular}
\end{center}
\caption{Overlap constants of plain $\HBs$ and of $\HBtp$; only the
right-hand column is used in this manuscript.}\label{s2:tabTauConstants}
\end{table}

\begin{remark}[constants of $\HBtp$ versus plain $\HBs$; \RTb{I8}]\label{s2:remTauConstants}
Plain $\HBs$ is the procedure of Definition~\ref{s2:defHBtp} in which the
second level of (R3) is the recursion of \cite[Lemma~14]{BM} at $s=s_l$ (no
$\tau$-rules) and (GC) is omitted; $\HBt$ has the $\tau$-rules but no (GC). The
overlap constants change because the separator of a $\tau$-split may exceed
the separator $N$ of \cite{BM} by $|\Hout|+|\Hin|$, which raises the charge in
Lemma~\ref{s2:lemOVgeneric} from $c=1$ to $c=1.6$ at the second level.
Table~\ref{s2:tabTauConstants} lists the values; only the right-hand column is
used in this manuscript. The left-hand column is recorded for comparison only;
it is neither used nor proved here.
\deps{\ref{s2:propOV}, \ref{s2:lem14tau}, \ref{s2:lemTower}, \ref{s2:defHBtp}, \ref{s2:lemOVgeneric}}
\fixes{\RTb{I8}}
\end{remark}

